% Môn: Vật lý
% Lớp: 11
% Chương: Chương II. Sóng
% Bài: L11C2 Bài 14. Bài tập về sóng
% Số câu: 174
% App đọc file này trực tiếp — sửa rồi Commit trên GitHub, bấm Đồng bộ GitHub .tex

% L11C2 Bài 14. Bài tập về sóng

% ===== Câu 1 =====
% ID: L11C2B14-01-DS
% Mức: TH
\dangbt{Bài toán giao thoa sóng}
\begin{ex}
% ID: L11C2B14-01-DS
% Mức: TH
Xét các phát biểu sau trong dạng “Bài toán giao thoa sóng”.
\choiceTF
{\True Một sóng có $v=30$ m/s, $f=5$ Hz. Bước sóng là — phát biểu đúng là: $6\,\text{m}$.}
{Một sóng có $v=30$ m/s, $f=5$ Hz. Bước sóng là — phát biểu $150\,\text{m}$ là đúng.}
{\True Độ lệch pha giữa hai điểm cách nhau $\Delta x$ trên phương truyền sóng là — phát biểu đúng là: $2\pi\Delta x/\lambda$.}
{Độ lệch pha giữa hai điểm cách nhau $\Delta x$ trên phương truyền sóng là — phát biểu $2\pi\lambda/\Delta x$ là đúng.}
\loigiai{
\begin{itemchoice}
\item (Đúng) Một sóng có $v=30$ m/s, $f=5$ Hz. Bước sóng là — phát biểu đúng là: $6\,\text{m}$. (Vì): Bước sóng được tính bằng công thức $\lambda = v/f$. Với $v=30$ m/s và $f=5$ Hz, ta có $\lambda = 30/5 = 6$ m
\item (Sai) Một sóng có $v=30$ m/s, $f=5$ Hz. Bước sóng là — phát biểu $150\,\text{m}$ là đúng. (Vì): Bước sóng được tính bằng công thức $\lambda = v/f$. Với $v=30$ m/s và $f=5$ Hz, ta có $\lambda = 30/5 = 6$ m, không phải $150\,\text{m}$
\item (Đúng) Độ lệch pha giữa hai điểm cách nhau $\Delta x$ trên phương truyền sóng là — phát biểu đúng là: $2\pi\Delta x/\lambda$. (Vì): ộ lệch pha giữa hai điểm trên phương truyền sóng cách nhau một khoảng $\Delta x$ được cho bởi công thức $|\Delta\varphi| = 2\pi\Delta x/\lambda$
\item (Sai) Độ lệch pha giữa hai điểm cách nhau $\Delta x$ trên phương truyền sóng là — phát biểu $2\pi\lambda/\Delta x$ là đúng. (Vì): Độ lệch pha giữa hai điểm trên phương truyền sóng cách nhau một khoảng $\Delta x$ được cho bởi công thức $|\Delta\varphi| = 2\pi\Delta x/\lambda$, không phải $2\pi\lambda/\Delta x$
\end{itemchoice}
}
\end{ex}

% ===== Câu 2 =====
% ID: L11C2B14-01-TL
% Mức: TH
\begin{bt}
% ID: L11C2B14-01-TL
% Mức: TH
Trong dạng “Bài toán giao thoa sóng”, Một sóng truyền với tốc độ $20\,\text{m}$ /s và chu kì $0,1\,\text{s}$. Tính bước sóng.
\loigiai{
$f=1/T=10$ Hz; $\lambda=v/f=20/10=2$ m.
}
\end{bt}

% ===== Câu 3 =====
% ID: L11C2B14-01-TLN
% Mức: TH
\begin{ex}
% ID: L11C2B14-01-TLN
% Mức: TH
Trong dạng “Bài toán giao thoa sóng”, Sóng có chu kì $0,25\,\text{s}$. Tính tần số theo Hz.
\shortans{4}
\loigiai{
$f=1/T=1/0,25=4$ Hz.
}
\end{ex}

% ===== Câu 4 =====
% ID: L11C2B14-01-TN
% Mức: NB
\begin{ex}
% ID: L11C2B14-01-TN
% Mức: NB
Trong dạng “Bài toán giao thoa sóng”, Một sóng có $v=30$ m/s, $f=5$ Hz. Bước sóng là
\choice
{\True $6\,\text{m}$}
{$150\,\text{m}$}
{$0,167\,\text{m}$}
{$35\,\text{m}$}
\loigiai{
Đáp án A. $\lambda=v/f=30/5=6$ m.
}
\end{ex}

% ===== Câu 5 =====
% ID: L11C2B14-02-DS
% Mức: TH
\begin{ex}
% ID: L11C2B14-02-DS
% Mức: TH
Xét các phát biểu sau trong dạng “Bài toán giao thoa sóng”.
\choiceTF
{\True Hai điểm cách nhau $\lambda/4$ trên phương truyền sóng lệch pha — phát biểu đúng là: $\pi/2$.}
{Hai điểm cách nhau $\lambda/4$ trên phương truyền sóng lệch pha — phát biểu $\pi$ là đúng.}
{\True Hai nguồn cùng pha, cực đại giao thoa thỏa — phát biểu đúng là: $\Delta d=k\lambda$.}
{Hai nguồn cùng pha, cực đại giao thoa thỏa — phát biểu $d_1+d_2=k\lambda$ là đúng.}
\loigiai{
\begin{itemchoice}
\item (Đúng) Hai điểm cách nhau $\lambda/4$ trên phương truyền sóng lệch pha — phát biểu đúng là: $\pi/2$. (Vì): $\Delta\varphi=2\pi(\lambda/4)/\lambda=\pi/2$. b) Sai: phương án nêu không phù hợp. c) Đúng: Điều kiện cực đại của hai nguồn cùng pha là $\Delta d=k\lambda$. d) Sai: phương án nêu không phù hợp
\item (Sai) Hai điểm cách nhau $\lambda/4$ trên phương truyền sóng lệch pha — phát biểu $\pi$ là đúng.
\item (Đúng) Hai nguồn cùng pha, cực đại giao thoa thỏa — phát biểu đúng là: $\Delta d=k\lambda$.
\item (Sai) Hai nguồn cùng pha, cực đại giao thoa thỏa — phát biểu $d_1+d_2=k\lambda$ là đúng.
\end{itemchoice}
}
\end{ex}

% ===== Câu 6 =====
% ID: L11C2B14-02-TL
% Mức: TH
\begin{ex}
% ID: L11C2B14-02-TL
% Mức: TH
Trong dạng “Bài toán giao thoa sóng”, Viết biểu thức độ lệch pha của hai điểm trên phương truyền sóng và nêu ý nghĩa.
\choiceTF

\loigiai{
$\Delta\varphi=2\pi\Delta x/\lambda$. Công thức cho biết quan hệ pha phụ thuộc khoảng cách theo phương truyền sóng.
}
\end{ex}

% ===== Câu 7 =====
% ID: L11C2B14-02-TLN
% Mức: TH
\begin{ex}
% ID: L11C2B14-02-TLN
% Mức: TH
Trong dạng “Bài toán giao thoa sóng”, Sóng có $v=24$ m/s và $f=8$ Hz. Tính bước sóng theo m.
\shortans{3}
\loigiai{
$\lambda=v/f=24/8=3$ m.
}
\end{ex}

% ===== Câu 8 =====
% ID: L11C2B14-02-TN
% Mức: NB
\begin{ex}
% ID: L11C2B14-02-TN
% Mức: NB
Trong dạng “Bài toán giao thoa sóng”, Độ lệch pha giữa hai điểm cách nhau $\Delta x$ trên phương truyền sóng là
\choice
{\True $2\pi\Delta x/\lambda$}
{$\pi\Delta x/\lambda$}
{$2\pi\lambda/\Delta x$}
{$\Delta x/(2\pi\lambda)$}
\loigiai{
Đáp án A. Theo phương trình sóng, $|\Delta\varphi|=2\pi\Delta x/\lambda$.
}
\end{ex}

% ===== Câu 9 =====
% ID: L11C2B14-03-DS
% Mức: VD
\begin{ex}
% ID: L11C2B14-03-DS
% Mức: VD
Xét các phát biểu sau trong dạng “Bài toán giao thoa sóng”.
\choiceTF
{\True Hai nút sóng dừng liên tiếp cách nhau — phát biểu đúng là: $\lambda/2$.}
{Hai nút sóng dừng liên tiếp cách nhau — phát biểu $\lambda$ là đúng.}
{\True Âm có tần số càng lớn thì cảm giác âm càng — phát biểu đúng là: cao.}
{Âm có tần số càng lớn thì cảm giác âm càng — phát biểu trầm là đúng.}
\loigiai{
\begin{itemchoice}
\item (Đúng) Hai nút sóng dừng liên tiếp cách nhau — phát biểu đúng là: $\lambda/2$. (Vì): Khoảng cách hai nút liên tiếp là nửa bước sóng. b) Sai: phương án nêu không phù hợp. c) Đúng: Độ cao của âm phụ thuộc chủ yếu vào tần số. d) Sai: phương án nêu không phù hợp
\item (Sai) Hai nút sóng dừng liên tiếp cách nhau — phát biểu $\lambda$ là đúng.
\item (Đúng) Âm có tần số càng lớn thì cảm giác âm càng — phát biểu đúng là: cao.
\item (Sai) Âm có tần số càng lớn thì cảm giác âm càng — phát biểu trầm là đúng.
\end{itemchoice}
}
\end{ex}

% ===== Câu 10 =====
% ID: L11C2B14-03-TL
% Mức: VD
\begin{ex}
% ID: L11C2B14-03-TL
% Mức: VD
Trong dạng “Bài toán giao thoa sóng”, Hai nguồn cùng pha có $\lambda=3$ cm; tại $M$, $d_1=14$ cm, $d_2=21,5$ cm. Xác định trạng thái giao thoa.
\choiceTF

\loigiai{
$\Delta d=7,5\,\mathrm{cm}=2,5\lambda$, nên $M$ là cực tiểu.
}
\end{ex}

% ===== Câu 11 =====
% ID: L11C2B14-03-TLN
% Mức: VD
\begin{ex}
% ID: L11C2B14-03-TLN
% Mức: VD
Trong dạng “Bài toán giao thoa sóng”, Hai điểm cách nhau $1\,\text{m}$ trên phương truyền sóng, $\lambda=4$ m. Tính độ lệch pha theo radian, lấy kết quả theo $\pi$.
\shortans{0,5}
\loigiai{
$\Delta\varphi=2\pi\Delta x/\lambda=2\pi/4=0,5\pi$ rad.
}
\end{ex}

% ===== Câu 12 =====
% ID: L11C2B14-03-TN
% Mức: TH
\begin{ex}
% ID: L11C2B14-03-TN
% Mức: TH
Trong dạng “Bài toán giao thoa sóng”, Hai điểm cách nhau $\lambda/4$ trên phương truyền sóng lệch pha
\choice
{\True $\pi/2$}
{$\pi$}
{$2\pi$}
{$\pi/4$}
\loigiai{
Đáp án A. $\Delta\varphi=2\pi(\lambda/4)/\lambda=\pi/2$.
}
\end{ex}

% ===== Câu 13 =====
% ID: L11C2B14-04-DS
% Mức: VD
\begin{ex}
% ID: L11C2B14-04-DS
% Mức: VD
Xét các phát biểu sau trong dạng “Bài toán giao thoa sóng”.
\choiceTF
{\True Mức cường độ âm được tính bởi — phát biểu đúng là: $L=10\log(I/I_0)$.}
{Mức cường độ âm được tính bởi — phát biểu $L=I/I_0$ là đúng.}
{\True Khi khoảng cách đến nguồn điểm tăng gấp đôi, cường độ âm trong môi trường không hấp thụ — phát biểu đúng là: giảm 4 lần.}
{Khi khoảng cách đến nguồn điểm tăng gấp đôi, cường độ âm trong môi trường không hấp thụ — phát biểu tăng 4 lần là đúng.}
\loigiai{
\begin{itemchoice}
\item (Đúng) Mức cường độ âm được tính bởi — phát biểu đúng là: $L=10\log(I/I_0)$. (Vì): Theo định nghĩa, $L=10\log_{10}(I/I_0)$, đơn vị dB. b) Sai: phương án nêu không phù hợp. c) Đúng: Với nguồn điểm $I\sim1/r^2$, nên r tăng 2 lần thì $I$ giảm 4 lần. d) Sai: phương án nêu không phù hợp
\item (Sai) Mức cường độ âm được tính bởi — phát biểu $L=I/I_0$ là đúng.
\item (Đúng) Khi khoảng cách đến nguồn điểm tăng gấp đôi, cường độ âm trong môi trường không hấp thụ — phát biểu đúng là: giảm 4 lần.
\item (Sai) Khi khoảng cách đến nguồn điểm tăng gấp đôi, cường độ âm trong môi trường không hấp thụ — phát biểu tăng 4 lần là đúng.
\end{itemchoice}
}
\end{ex}

% ===== Câu 14 =====
% ID: L11C2B14-04-TL
% Mức: VD
\begin{bt}
% ID: L11C2B14-04-TL
% Mức: VD
Trong dạng “Bài toán giao thoa sóng”, Dây hai đầu cố định dài $1,2\,\text{m}$ có tốc độ sóng $48\,\text{m}$ /s. Tính tần số họa âm thứ ba.
\loigiai{
$f_1=v/(2L)=48/2,4=20$ Hz; $f_3=3f_1=60$ Hz.
}
\end{bt}

% ===== Câu 15 =====
% ID: L11C2B14-04-TLN
% Mức: VD
\begin{ex}
% ID: L11C2B14-04-TLN
% Mức: VD
Trong dạng “Bài toán giao thoa sóng”, Cường độ âm tăng 100 lần. Mức cường độ âm tăng bao nhiêu dB?
\shortans{20}
\loigiai{
$\Delta L=10\log100=20$ dB.
}
\end{ex}

% ===== Câu 16 =====
% ID: L11C2B14-04-TN
% Mức: TH
\begin{ex}
% ID: L11C2B14-04-TN
% Mức: TH
Trong dạng “Bài toán giao thoa sóng”, Hai nguồn cùng pha, cực đại giao thoa thỏa
\choice
{\True $\Delta d=k\lambda$}
{$\Delta d=(k+1/2)\lambda$}
{$d_1+d_2=k\lambda$}
{$\Delta d=\lambda/4$}
\loigiai{
Đáp án A. Điều kiện cực đại của hai nguồn cùng pha là $\Delta d=k\lambda$.
}
\end{ex}

% ===== Câu 17 =====
% ID: L11C2B14-05-DS
% Mức: VDC
\begin{ex}
% ID: L11C2B14-05-DS
% Mức: VDC
Xét các phát biểu sau trong dạng “Bài toán giao thoa sóng”.
\choiceTF
{\True Sóng truyền từ môi trường này sang môi trường khác, đại lượng không đổi là — phát biểu đúng là: tần số.}
{Sóng truyền từ môi trường này sang môi trường khác, đại lượng không đổi là — phát biểu bước sóng là đúng.}
{\True Một dây hai đầu cố định dài $0,6\,\text{m}$ có 3 bụng. Bước sóng là — phát biểu đúng là: $0,4\,\text{m}$.}
{Một dây hai đầu cố định dài $0,6\,\text{m}$ có 3 bụng. Bước sóng là — phát biểu $1,8\,\text{m}$ là đúng.}
\loigiai{
\begin{itemchoice}
\item (Đúng) Sóng truyền từ môi trường này sang môi trường khác, đại lượng không đổi là — phát biểu đúng là: tần số. (Vì): Tần số do nguồn quyết định và không đổi khi qua mặt phân cách. b) Sai: phương án nêu không phù hợp. c) Đúng: $0,6=3\lambda/2$, suy ra $\lambda=0,4$ m. d) Sai: phương án nêu không phù hợp
\item (Sai) Sóng truyền từ môi trường này sang môi trường khác, đại lượng không đổi là — phát biểu bước sóng là đúng.
\item (Đúng) Một dây hai đầu cố định dài $0,6\,\text{m}$ có 3 bụng. Bước sóng là — phát biểu đúng là: $0,4\,\text{m}$.
\item (Sai) Một dây hai đầu cố định dài $0,6\,\text{m}$ có 3 bụng. Bước sóng là — phát biểu $1,8\,\text{m}$ là đúng.
\end{itemchoice}
}
\end{ex}

% ===== Câu 18 =====
% ID: L11C2B14-05-TL
% Mức: VD
\begin{ex}
% ID: L11C2B14-05-TL
% Mức: VD
Trong dạng “Bài toán giao thoa sóng”, Giải thích sự thay đổi bước sóng khi sóng truyền sang môi trường khác.
\choiceTF

\loigiai{
Tần số do nguồn quyết định nên không đổi; tốc độ phụ thuộc môi trường, vì $\lambda=v/f$ nên bước sóng thay đổi theo tốc độ.
}
\end{ex}

% ===== Câu 19 =====
% ID: L11C2B14-05-TLN
% Mức: VD
\begin{ex}
% ID: L11C2B14-05-TLN
% Mức: VD
Trong dạng “Bài toán giao thoa sóng”, Dây hai đầu cố định dài $1\,\text{m}$ có bước sóng $0,5\,\text{m}$. Số bụng là bao nhiêu?
\shortans{4}
\loigiai{
$k=2L/\lambda=2/0,5=4$ bụng.
}
\end{ex}

% ===== Câu 20 =====
% ID: L11C2B14-05-TN
% Mức: TH
\begin{ex}
% ID: L11C2B14-05-TN
% Mức: TH
Trong dạng “Bài toán giao thoa sóng”, Hai nút sóng dừng liên tiếp cách nhau
\choice
{\True $\lambda/2$}
{$\lambda$}
{$\lambda/4$}
{$2\lambda$}
\loigiai{
Đáp án A. Khoảng cách hai nút liên tiếp là nửa bước sóng.
}
\end{ex}

% ===== Câu 21 =====
% ID: L11C2B14-06-DS
% Mức: TH
\dangbt{Bài toán sóng dừng}
\begin{ex}
% ID: L11C2B14-06-DS
% Mức: TH
Xét các phát biểu sau trong dạng “Bài toán sóng dừng”.
\choiceTF
{\True Một sóng có $v=30$ m/s, $f=5$ Hz. Bước sóng là — phát biểu đúng là: $6\,\text{m}$.}
{Một sóng có $v=30$ m/s, $f=5$ Hz. Bước sóng là — phát biểu $150\,\text{m}$ là đúng.}
{\True Độ lệch pha giữa hai điểm cách nhau $\Delta x$ trên phương truyền sóng là — phát biểu đúng là: $2\pi\Delta x/\lambda$.}
{Độ lệch pha giữa hai điểm cách nhau $\Delta x$ trên phương truyền sóng là — phát biểu $2\pi\lambda/\Delta x$ là đúng.}
\loigiai{
\begin{itemchoice}
\item (Đúng) Một sóng có $v=30$ m/s, $f=5$ Hz. Bước sóng là — phát biểu đúng là: $6\,\text{m}$. (Vì): $\lambda=v/f=30/5=6$ m. b) Sai: phương án nêu không phù hợp. c) Đúng: Theo phương trình sóng, $|\Delta\varphi|=2\pi\Delta x/\lambda$. d) Sai: phương án nêu không phù hợp
\item (Sai) Một sóng có $v=30$ m/s, $f=5$ Hz. Bước sóng là — phát biểu $150\,\text{m}$ là đúng.
\item (Đúng) Độ lệch pha giữa hai điểm cách nhau $\Delta x$ trên phương truyền sóng là — phát biểu đúng là: $2\pi\Delta x/\lambda$.
\item (Sai) Độ lệch pha giữa hai điểm cách nhau $\Delta x$ trên phương truyền sóng là — phát biểu $2\pi\lambda/\Delta x$ là đúng.
\end{itemchoice}
}
\end{ex}

% ===== Câu 22 =====
% ID: L11C2B14-06-TL
% Mức: VDC
\dangbt{Bài toán giao thoa sóng}
\begin{ex}
% ID: L11C2B14-06-TL
% Mức: VDC
Trong dạng “Bài toán giao thoa sóng”, Tại sao đứng xa nguồn âm thì âm nghe nhỏ hơn?
\choiceTF

\loigiai{
Năng lượng sóng phân bố trên diện tích ngày càng lớn; với nguồn điểm, cường độ giảm theo $1/r^2$, nên âm nghe nhỏ hơn.
}
\end{ex}

% ===== Câu 23 =====
% ID: L11C2B14-06-TLN
% Mức: VDC
\begin{ex}
% ID: L11C2B14-06-TLN
% Mức: VDC
Trong dạng “Bài toán giao thoa sóng”, Nguồn điểm có cường độ âm 8 đơn vị tại r. Ở khoảng cách 2r, cường độ theo cùng đơn vị là bao nhiêu?
\shortans{2}
\loigiai{
$I_2=I_1(r/2r)^2=8/4=2$.
}
\end{ex}

% ===== Câu 24 =====
% ID: L11C2B14-06-TN
% Mức: VD
\begin{ex}
% ID: L11C2B14-06-TN
% Mức: VD
Trong dạng “Bài toán giao thoa sóng”, Âm có tần số càng lớn thì cảm giác âm càng
\choice
{\True cao}
{to}
{trầm}
{nhỏ}
\loigiai{
Đáp án A. Độ cao của âm phụ thuộc chủ yếu vào tần số.
}
\end{ex}

% ===== Câu 25 =====
% ID: L11C2B14-07-DS
% Mức: TH
\dangbt{Bài toán sóng dừng}
\begin{ex}
% ID: L11C2B14-07-DS
% Mức: TH
Xét các phát biểu sau trong dạng “Bài toán sóng dừng”.
\choiceTF
{\True Hai điểm cách nhau $\lambda/4$ trên phương truyền sóng lệch pha — phát biểu đúng là: $\pi/2$.}
{Hai điểm cách nhau $\lambda/4$ trên phương truyền sóng lệch pha — phát biểu $\pi$ là đúng.}
{\True Hai nguồn cùng pha, cực đại giao thoa thỏa — phát biểu đúng là: $\Delta d=k\lambda$.}
{Hai nguồn cùng pha, cực đại giao thoa thỏa — phát biểu $d_1+d_2=k\lambda$ là đúng.}
\loigiai{
\begin{itemchoice}
\item (Đúng) Hai điểm cách nhau $\lambda/4$ trên phương truyền sóng lệch pha — phát biểu đúng là: $\pi/2$. (Vì): $\Delta\varphi=2\pi(\lambda/4)/\lambda=\pi/2$. b) Sai: phương án nêu không phù hợp. c) Đúng: Điều kiện cực đại của hai nguồn cùng pha là $\Delta d=k\lambda$. d) Sai: phương án nêu không phù hợp
\item (Sai) Hai điểm cách nhau $\lambda/4$ trên phương truyền sóng lệch pha — phát biểu $\pi$ là đúng.
\item (Đúng) Hai nguồn cùng pha, cực đại giao thoa thỏa — phát biểu đúng là: $\Delta d=k\lambda$.
\item (Sai) Hai nguồn cùng pha, cực đại giao thoa thỏa — phát biểu $d_1+d_2=k\lambda$ là đúng.
\end{itemchoice}
}
\end{ex}

% ===== Câu 26 =====
% ID: L11C2B14-07-TL
% Mức: TH
\begin{bt}
% ID: L11C2B14-07-TL
% Mức: TH
Trong dạng “Bài toán sóng dừng”, Một sóng truyền với tốc độ $20\,\text{m}$ /s và chu kì $0,1\,\text{s}$. Tính bước sóng.
\loigiai{
$f=1/T=10$ Hz; $\lambda=v/f=20/10=2$ m.
}
\end{bt}

% ===== Câu 27 =====
% ID: L11C2B14-07-TLN
% Mức: TH
\begin{ex}
% ID: L11C2B14-07-TLN
% Mức: TH
Trong dạng “Bài toán sóng dừng”, Sóng có chu kì $0,25\,\text{s}$. Tính tần số theo Hz.
\shortans{4}
\loigiai{
$f=1/T=1/0,25=4$ Hz.
}
\end{ex}

% ===== Câu 28 =====
% ID: L11C2B14-07-TN
% Mức: VD
\dangbt{Bài toán giao thoa sóng}
\begin{ex}
% ID: L11C2B14-07-TN
% Mức: VD
Trong dạng “Bài toán giao thoa sóng”, Mức cường độ âm được tính bởi
\choice
{\True $L=10\log(I/I_0)$}
{$L=I/I_0$}
{$L=20I/I_0$}
{$L=10I_0/I$}
\loigiai{
Đáp án A. Theo định nghĩa, $L=10\log_{10}(I/I_0)$, đơn vị dB.
}
\end{ex}

% ===== Câu 29 =====
% ID: L11C2B14-08-DS
% Mức: VD
\dangbt{Bài toán sóng dừng}
\begin{ex}
% ID: L11C2B14-08-DS
% Mức: VD
Xét các phát biểu sau trong dạng “Bài toán sóng dừng”.
\choiceTF
{\True Hai nút sóng dừng liên tiếp cách nhau — phát biểu đúng là: $\lambda/2$.}
{Hai nút sóng dừng liên tiếp cách nhau — phát biểu $\lambda$ là đúng.}
{\True Âm có tần số càng lớn thì cảm giác âm càng — phát biểu đúng là: cao.}
{Âm có tần số càng lớn thì cảm giác âm càng — phát biểu trầm là đúng.}
\loigiai{
\begin{itemchoice}
\item (Đúng) Hai nút sóng dừng liên tiếp cách nhau — phát biểu đúng là: $\lambda/2$. (Vì): Khoảng cách hai nút liên tiếp là nửa bước sóng. b) Sai: phương án nêu không phù hợp. c) Đúng: Độ cao của âm phụ thuộc chủ yếu vào tần số. d) Sai: phương án nêu không phù hợp
\item (Sai) Hai nút sóng dừng liên tiếp cách nhau — phát biểu $\lambda$ là đúng.
\item (Đúng) Âm có tần số càng lớn thì cảm giác âm càng — phát biểu đúng là: cao.
\item (Sai) Âm có tần số càng lớn thì cảm giác âm càng — phát biểu trầm là đúng.
\end{itemchoice}
}
\end{ex}

% ===== Câu 30 =====
% ID: L11C2B14-08-TL
% Mức: TH
\begin{ex}
% ID: L11C2B14-08-TL
% Mức: TH
Trong dạng “Bài toán sóng dừng”, Viết biểu thức độ lệch pha của hai điểm trên phương truyền sóng và nêu ý nghĩa.
\choiceTF

\loigiai{
$\Delta\varphi=2\pi\Delta x/\lambda$. Công thức cho biết quan hệ pha phụ thuộc khoảng cách theo phương truyền sóng.
}
\end{ex}

% ===== Câu 31 =====
% ID: L11C2B14-08-TLN
% Mức: TH
\begin{ex}
% ID: L11C2B14-08-TLN
% Mức: TH
Trong dạng “Bài toán sóng dừng”, Sóng có $v=24$ m/s và $f=8$ Hz. Tính bước sóng theo m.
\shortans{3}
\loigiai{
$\lambda=v/f=24/8=3$ m.
}
\end{ex}

% ===== Câu 32 =====
% ID: L11C2B14-08-TN
% Mức: VD
\dangbt{Bài toán giao thoa sóng}
\begin{ex}
% ID: L11C2B14-08-TN
% Mức: VD
Trong dạng “Bài toán giao thoa sóng”, Khi khoảng cách đến nguồn điểm tăng gấp đôi, cường độ âm trong môi trường không hấp thụ
\choice
{\True giảm 4 lần}
{giảm 2 lần}
{tăng 4 lần}
{không đổi}
\loigiai{
Đáp án A. Với nguồn điểm $I\sim1/r^2$, nên r tăng 2 lần thì $I$ giảm 4 lần.
}
\end{ex}

% ===== Câu 33 =====
% ID: L11C2B14-09-DS
% Mức: VD
\dangbt{Bài toán sóng dừng}
\begin{ex}
% ID: L11C2B14-09-DS
% Mức: VD
Xét các phát biểu sau trong dạng “Bài toán sóng dừng”.
\choiceTF
{\True Mức cường độ âm được tính bởi — phát biểu đúng là: $L=10\log(I/I_0)$.}
{Mức cường độ âm được tính bởi — phát biểu $L=I/I_0$ là đúng.}
{\True Khi khoảng cách đến nguồn điểm tăng gấp đôi, cường độ âm trong môi trường không hấp thụ — phát biểu đúng là: giảm 4 lần.}
{Khi khoảng cách đến nguồn điểm tăng gấp đôi, cường độ âm trong môi trường không hấp thụ — phát biểu tăng 4 lần là đúng.}
\loigiai{
\begin{itemchoice}
\item (Đúng) Mức cường độ âm được tính bởi — phát biểu đúng là: $L=10\log(I/I_0)$. (Vì): Theo định nghĩa, $L=10\log_{10}(I/I_0)$, đơn vị dB. b) Sai: phương án nêu không phù hợp. c) Đúng: Với nguồn điểm $I\sim1/r^2$, nên r tăng 2 lần thì $I$ giảm 4 lần. d) Sai: phương án nêu không phù hợp
\item (Sai) Mức cường độ âm được tính bởi — phát biểu $L=I/I_0$ là đúng.
\item (Đúng) Khi khoảng cách đến nguồn điểm tăng gấp đôi, cường độ âm trong môi trường không hấp thụ — phát biểu đúng là: giảm 4 lần.
\item (Sai) Khi khoảng cách đến nguồn điểm tăng gấp đôi, cường độ âm trong môi trường không hấp thụ — phát biểu tăng 4 lần là đúng.
\end{itemchoice}
}
\end{ex}

% ===== Câu 34 =====
% ID: L11C2B14-09-TL
% Mức: VD
\begin{ex}
% ID: L11C2B14-09-TL
% Mức: VD
Trong dạng “Bài toán sóng dừng”, Hai nguồn cùng pha có $\lambda=3$ cm; tại $M$, $d_1=14$ cm, $d_2=21,5$ cm. Xác định trạng thái giao thoa.
\choiceTF

\loigiai{
$\Delta d=7,5\,\mathrm{cm}=2,5\lambda$, nên $M$ là cực tiểu.
}
\end{ex}

% ===== Câu 35 =====
% ID: L11C2B14-09-TLN
% Mức: VD
\begin{ex}
% ID: L11C2B14-09-TLN
% Mức: VD
Trong dạng “Bài toán sóng dừng”, Hai điểm cách nhau $1\,\text{m}$ trên phương truyền sóng, $\lambda=4$ m. Tính độ lệch pha theo radian, lấy kết quả theo $\pi$.
\shortans{0,5}
\loigiai{
$\Delta\varphi=2\pi\Delta x/\lambda=2\pi/4=0,5\pi$ rad.
}
\end{ex}

% ===== Câu 36 =====
% ID: L11C2B14-09-TN
% Mức: VDC
\dangbt{Bài toán giao thoa sóng}
\begin{ex}
% ID: L11C2B14-09-TN
% Mức: VDC
Trong dạng “Bài toán giao thoa sóng”, Sóng truyền từ môi trường này sang môi trường khác, đại lượng không đổi là
\choice
{\True tần số}
{bước sóng}
{tốc độ}
{biên độ}
\loigiai{
Đáp án A. Tần số do nguồn quyết định và không đổi khi qua mặt phân cách.
}
\end{ex}

% ===== Câu 37 =====
% ID: L11C2B14-10-DS
% Mức: VDC
\dangbt{Bài toán sóng dừng}
\begin{ex}
% ID: L11C2B14-10-DS
% Mức: VDC
Xét các phát biểu sau trong dạng “Bài toán sóng dừng”.
\choiceTF
{\True Sóng truyền từ môi trường này sang môi trường khác, đại lượng không đổi là — phát biểu đúng là: tần số.}
{Sóng truyền từ môi trường này sang môi trường khác, đại lượng không đổi là — phát biểu bước sóng là đúng.}
{\True Một dây hai đầu cố định dài $0,6\,\text{m}$ có 3 bụng. Bước sóng là — phát biểu đúng là: $0,4\,\text{m}$.}
{Một dây hai đầu cố định dài $0,6\,\text{m}$ có 3 bụng. Bước sóng là — phát biểu $1,8\,\text{m}$ là đúng.}
\loigiai{
\begin{itemchoice}
\item (Đúng) Sóng truyền từ môi trường này sang môi trường khác, đại lượng không đổi là — phát biểu đúng là: tần số. (Vì): Tần số do nguồn quyết định và không đổi khi qua mặt phân cách. b) Sai: phương án nêu không phù hợp. c) Đúng: $0,6=3\lambda/2$, suy ra $\lambda=0,4$ m. d) Sai: phương án nêu không phù hợp
\item (Sai) Sóng truyền từ môi trường này sang môi trường khác, đại lượng không đổi là — phát biểu bước sóng là đúng.
\item (Đúng) Một dây hai đầu cố định dài $0,6\,\text{m}$ có 3 bụng. Bước sóng là — phát biểu đúng là: $0,4\,\text{m}$.
\item (Sai) Một dây hai đầu cố định dài $0,6\,\text{m}$ có 3 bụng. Bước sóng là — phát biểu $1,8\,\text{m}$ là đúng.
\end{itemchoice}
}
\end{ex}

% ===== Câu 38 =====
% ID: L11C2B14-10-TL
% Mức: VD
\begin{bt}
% ID: L11C2B14-10-TL
% Mức: VD
Trong dạng “Bài toán sóng dừng”, Dây hai đầu cố định dài $1,2\,\text{m}$ có tốc độ sóng $48\,\text{m}$ /s. Tính tần số họa âm thứ ba.
\loigiai{
$f_1=v/(2L)=48/2,4=20$ Hz; $f_3=3f_1=60$ Hz.
}
\end{bt}

% ===== Câu 39 =====
% ID: L11C2B14-10-TLN
% Mức: VD
\begin{ex}
% ID: L11C2B14-10-TLN
% Mức: VD
Trong dạng “Bài toán sóng dừng”, Cường độ âm tăng 100 lần. Mức cường độ âm tăng bao nhiêu dB?
\shortans{20}
\loigiai{
$\Delta L=10\log100=20$ dB.
}
\end{ex}

% ===== Câu 40 =====
% ID: L11C2B14-10-TN
% Mức: VDC
\dangbt{Bài toán giao thoa sóng}
\begin{ex}
% ID: L11C2B14-10-TN
% Mức: VDC
Trong dạng “Bài toán giao thoa sóng”, Một dây hai đầu cố định dài $0,6\,\text{m}$ có 3 bụng. Bước sóng là
\choice
{\True $0,4\,\text{m}$}
{$0,2\,\text{m}$}
{$1,8\,\text{m}$}
{$0,6\,\text{m}$}
\loigiai{
Đáp án A. $0,6=3\lambda/2$, suy ra $\lambda=0,4$ m.
}
\end{ex}

% ===== Câu 41 =====
% ID: L11C2B14-11-DS
% Mức: TH
\dangbt{Bài toán âm, tần số và cường độ âm}
\begin{ex}
% ID: L11C2B14-11-DS
% Mức: TH
Xét các phát biểu sau trong dạng “Bài toán âm, tần số và cường độ âm”.
\choiceTF
{\True Một sóng có $v=30$ m/s, $f=5$ Hz. Bước sóng là — phát biểu đúng là: $6\,\text{m}$.}
{Một sóng có $v=30$ m/s, $f=5$ Hz. Bước sóng là — phát biểu $150\,\text{m}$ là đúng.}
{\True Độ lệch pha giữa hai điểm cách nhau $\Delta x$ trên phương truyền sóng là — phát biểu đúng là: $2\pi\Delta x/\lambda$.}
{Độ lệch pha giữa hai điểm cách nhau $\Delta x$ trên phương truyền sóng là — phát biểu $2\pi\lambda/\Delta x$ là đúng.}
\loigiai{
\begin{itemchoice}
\item (Đúng) Một sóng có $v=30$ m/s, $f=5$ Hz. Bước sóng là — phát biểu đúng là: $6\,\text{m}$. (Vì): $\lambda=v/f=30/5=6$ m. b) Sai: phương án nêu không phù hợp. c) Đúng: Theo phương trình sóng, $|\Delta\varphi|=2\pi\Delta x/\lambda$. d) Sai: phương án nêu không phù hợp
\item (Sai) Một sóng có $v=30$ m/s, $f=5$ Hz. Bước sóng là — phát biểu $150\,\text{m}$ là đúng.
\item (Đúng) Độ lệch pha giữa hai điểm cách nhau $\Delta x$ trên phương truyền sóng là — phát biểu đúng là: $2\pi\Delta x/\lambda$.
\item (Sai) Độ lệch pha giữa hai điểm cách nhau $\Delta x$ trên phương truyền sóng là — phát biểu $2\pi\lambda/\Delta x$ là đúng.
\end{itemchoice}
}
\end{ex}

% ===== Câu 42 =====
% ID: L11C2B14-11-TL
% Mức: VD
\dangbt{Bài toán sóng dừng}
\begin{ex}
% ID: L11C2B14-11-TL
% Mức: VD
Trong dạng “Bài toán sóng dừng”, Giải thích sự thay đổi bước sóng khi sóng truyền sang môi trường khác.
\choiceTF

\loigiai{
Tần số do nguồn quyết định nên không đổi; tốc độ phụ thuộc môi trường, vì $\lambda=v/f$ nên bước sóng thay đổi theo tốc độ.
}
\end{ex}

% ===== Câu 43 =====
% ID: L11C2B14-11-TLN
% Mức: VD
\begin{ex}
% ID: L11C2B14-11-TLN
% Mức: VD
Trong dạng “Bài toán sóng dừng”, Dây hai đầu cố định dài $1\,\text{m}$ có bước sóng $0,5\,\text{m}$. Số bụng là bao nhiêu?
\shortans{4}
\loigiai{
$k=2L/\lambda=2/0,5=4$ bụng.
}
\end{ex}

% ===== Câu 44 =====
% ID: L11C2B14-11-TN
% Mức: NB
\begin{ex}
% ID: L11C2B14-11-TN
% Mức: NB
Trong dạng “Bài toán sóng dừng”, Một sóng có $v=30$ m/s, $f=5$ Hz. Bước sóng là
\choice
{\True $6\,\text{m}$}
{$150\,\text{m}$}
{$0,167\,\text{m}$}
{$35\,\text{m}$}
\loigiai{
Đáp án A. $\lambda=v/f=30/5=6$ m.
}
\end{ex}

% ===== Câu 45 =====
% ID: L11C2B14-12-DS
% Mức: TH
\dangbt{Bài toán âm, tần số và cường độ âm}
\begin{ex}
% ID: L11C2B14-12-DS
% Mức: TH
Xét các phát biểu sau trong dạng “Bài toán âm, tần số và cường độ âm”.
\choiceTF
{\True Hai điểm cách nhau $\lambda/4$ trên phương truyền sóng lệch pha — phát biểu đúng là: $\pi/2$.}
{Hai điểm cách nhau $\lambda/4$ trên phương truyền sóng lệch pha — phát biểu $\pi$ là đúng.}
{\True Hai nguồn cùng pha, cực đại giao thoa thỏa — phát biểu đúng là: $\Delta d=k\lambda$.}
{Hai nguồn cùng pha, cực đại giao thoa thỏa — phát biểu $d_1+d_2=k\lambda$ là đúng.}
\loigiai{
\begin{itemchoice}
\item (Đúng) Hai điểm cách nhau $\lambda/4$ trên phương truyền sóng lệch pha — phát biểu đúng là: $\pi/2$. (Vì): $\Delta\varphi=2\pi(\lambda/4)/\lambda=\pi/2$. b) Sai: phương án nêu không phù hợp. c) Đúng: Điều kiện cực đại của hai nguồn cùng pha là $\Delta d=k\lambda$. d) Sai: phương án nêu không phù hợp
\item (Sai) Hai điểm cách nhau $\lambda/4$ trên phương truyền sóng lệch pha — phát biểu $\pi$ là đúng.
\item (Đúng) Hai nguồn cùng pha, cực đại giao thoa thỏa — phát biểu đúng là: $\Delta d=k\lambda$.
\item (Sai) Hai nguồn cùng pha, cực đại giao thoa thỏa — phát biểu $d_1+d_2=k\lambda$ là đúng.
\end{itemchoice}
}
\end{ex}

% ===== Câu 46 =====
% ID: L11C2B14-12-TL
% Mức: VDC
\dangbt{Bài toán sóng dừng}
\begin{ex}
% ID: L11C2B14-12-TL
% Mức: VDC
Trong dạng “Bài toán sóng dừng”, Tại sao đứng xa nguồn âm thì âm nghe nhỏ hơn?
\choiceTF

\loigiai{
Năng lượng sóng phân bố trên diện tích ngày càng lớn; với nguồn điểm, cường độ giảm theo $1/r^2$, nên âm nghe nhỏ hơn.
}
\end{ex}

% ===== Câu 47 =====
% ID: L11C2B14-12-TLN
% Mức: VDC
\begin{ex}
% ID: L11C2B14-12-TLN
% Mức: VDC
Trong dạng “Bài toán sóng dừng”, Nguồn điểm có cường độ âm 8 đơn vị tại r. Ở khoảng cách 2r, cường độ theo cùng đơn vị là bao nhiêu?
\shortans{2}
\loigiai{
$I_2=I_1(r/2r)^2=8/4=2$.
}
\end{ex}

% ===== Câu 48 =====
% ID: L11C2B14-12-TN
% Mức: NB
\begin{ex}
% ID: L11C2B14-12-TN
% Mức: NB
Trong dạng “Bài toán sóng dừng”, Độ lệch pha giữa hai điểm cách nhau $\Delta x$ trên phương truyền sóng là
\choice
{\True $2\pi\Delta x/\lambda$}
{$\pi\Delta x/\lambda$}
{$2\pi\lambda/\Delta x$}
{$\Delta x/(2\pi\lambda)$}
\loigiai{
Đáp án A. Theo phương trình sóng, $|\Delta\varphi|=2\pi\Delta x/\lambda$.
}
\end{ex}

% ===== Câu 49 =====
% ID: L11C2B14-13-DS
% Mức: VD
\dangbt{Bài toán âm, tần số và cường độ âm}
\begin{ex}
% ID: L11C2B14-13-DS
% Mức: VD
Xét các phát biểu sau trong dạng “Bài toán âm, tần số và cường độ âm”.
\choiceTF
{\True Hai nút sóng dừng liên tiếp cách nhau — phát biểu đúng là: $\lambda/2$.}
{Hai nút sóng dừng liên tiếp cách nhau — phát biểu $\lambda$ là đúng.}
{\True Âm có tần số càng lớn thì cảm giác âm càng — phát biểu đúng là: cao.}
{Âm có tần số càng lớn thì cảm giác âm càng — phát biểu trầm là đúng.}
\loigiai{
\begin{itemchoice}
\item (Đúng) Hai nút sóng dừng liên tiếp cách nhau — phát biểu đúng là: $\lambda/2$. (Vì): Khoảng cách hai nút liên tiếp là nửa bước sóng. b) Sai: phương án nêu không phù hợp. c) Đúng: Độ cao của âm phụ thuộc chủ yếu vào tần số. d) Sai: phương án nêu không phù hợp
\item (Sai) Hai nút sóng dừng liên tiếp cách nhau — phát biểu $\lambda$ là đúng.
\item (Đúng) Âm có tần số càng lớn thì cảm giác âm càng — phát biểu đúng là: cao.
\item (Sai) Âm có tần số càng lớn thì cảm giác âm càng — phát biểu trầm là đúng.
\end{itemchoice}
}
\end{ex}

% ===== Câu 50 =====
% ID: L11C2B14-13-TL
% Mức: TH
\begin{bt}
% ID: L11C2B14-13-TL
% Mức: TH
Trong dạng “Bài toán âm, tần số và cường độ âm”, Một sóng truyền với tốc độ $20\,\text{m}$ /s và chu kì $0,1\,\text{s}$. Tính bước sóng.
\loigiai{
$f=1/T=10$ Hz; $\lambda=v/f=20/10=2$ m.
}
\end{bt}

% ===== Câu 51 =====
% ID: L11C2B14-13-TLN
% Mức: TH
\begin{ex}
% ID: L11C2B14-13-TLN
% Mức: TH
Trong dạng “Bài toán âm, tần số và cường độ âm”, Sóng có chu kì $0,25\,\text{s}$. Tính tần số theo Hz.
\shortans{4}
\loigiai{
$f=1/T=1/0,25=4$ Hz.
}
\end{ex}

% ===== Câu 52 =====
% ID: L11C2B14-13-TN
% Mức: TH
\dangbt{Bài toán sóng dừng}
\begin{ex}
% ID: L11C2B14-13-TN
% Mức: TH
Trong dạng “Bài toán sóng dừng”, Hai điểm cách nhau $\lambda/4$ trên phương truyền sóng lệch pha
\choice
{\True $\pi/2$}
{$\pi$}
{$2\pi$}
{$\pi/4$}
\loigiai{
Đáp án A. $\Delta\varphi=2\pi(\lambda/4)/\lambda=\pi/2$.
}
\end{ex}

% ===== Câu 53 =====
% ID: L11C2B14-14-DS
% Mức: VD
\dangbt{Bài toán âm, tần số và cường độ âm}
\begin{ex}
% ID: L11C2B14-14-DS
% Mức: VD
Xét các phát biểu sau trong dạng “Bài toán âm, tần số và cường độ âm”.
\choiceTF
{\True Mức cường độ âm được tính bởi — phát biểu đúng là: $L=10\log(I/I_0)$.}
{Mức cường độ âm được tính bởi — phát biểu $L=I/I_0$ là đúng.}
{\True Khi khoảng cách đến nguồn điểm tăng gấp đôi, cường độ âm trong môi trường không hấp thụ — phát biểu đúng là: giảm 4 lần.}
{Khi khoảng cách đến nguồn điểm tăng gấp đôi, cường độ âm trong môi trường không hấp thụ — phát biểu tăng 4 lần là đúng.}
\loigiai{
\begin{itemchoice}
\item (Đúng) Mức cường độ âm được tính bởi — phát biểu đúng là: $L=10\log(I/I_0)$. (Vì): Theo định nghĩa, $L=10\log_{10}(I/I_0)$, đơn vị dB. b) Sai: phương án nêu không phù hợp. c) Đúng: Với nguồn điểm $I\sim1/r^2$, nên r tăng 2 lần thì $I$ giảm 4 lần. d) Sai: phương án nêu không phù hợp
\item (Sai) Mức cường độ âm được tính bởi — phát biểu $L=I/I_0$ là đúng.
\item (Đúng) Khi khoảng cách đến nguồn điểm tăng gấp đôi, cường độ âm trong môi trường không hấp thụ — phát biểu đúng là: giảm 4 lần.
\item (Sai) Khi khoảng cách đến nguồn điểm tăng gấp đôi, cường độ âm trong môi trường không hấp thụ — phát biểu tăng 4 lần là đúng.
\end{itemchoice}
}
\end{ex}

% ===== Câu 54 =====
% ID: L11C2B14-14-TL
% Mức: TH
\begin{ex}
% ID: L11C2B14-14-TL
% Mức: TH
Trong dạng “Bài toán âm, tần số và cường độ âm”, Viết biểu thức độ lệch pha của hai điểm trên phương truyền sóng và nêu ý nghĩa.
\choiceTF

\loigiai{
$\Delta\varphi=2\pi\Delta x/\lambda$. Công thức cho biết quan hệ pha phụ thuộc khoảng cách theo phương truyền sóng.
}
\end{ex}

% ===== Câu 55 =====
% ID: L11C2B14-14-TLN
% Mức: TH
\begin{ex}
% ID: L11C2B14-14-TLN
% Mức: TH
Trong dạng “Bài toán âm, tần số và cường độ âm”, Sóng có $v=24$ m/s và $f=8$ Hz. Tính bước sóng theo m.
\shortans{3}
\loigiai{
$\lambda=v/f=24/8=3$ m.
}
\end{ex}

% ===== Câu 56 =====
% ID: L11C2B14-14-TN
% Mức: TH
\dangbt{Bài toán sóng dừng}
\begin{ex}
% ID: L11C2B14-14-TN
% Mức: TH
Trong dạng “Bài toán sóng dừng”, Hai nguồn cùng pha, cực đại giao thoa thỏa
\choice
{\True $\Delta d=k\lambda$}
{$\Delta d=(k+1/2)\lambda$}
{$d_1+d_2=k\lambda$}
{$\Delta d=\lambda/4$}
\loigiai{
Đáp án A. Điều kiện cực đại của hai nguồn cùng pha là $\Delta d=k\lambda$.
}
\end{ex}

% ===== Câu 57 =====
% ID: L11C2B14-15-DS
% Mức: VDC
\dangbt{Bài toán âm, tần số và cường độ âm}
\begin{ex}
% ID: L11C2B14-15-DS
% Mức: VDC
Xét các phát biểu sau trong dạng “Bài toán âm, tần số và cường độ âm”.
\choiceTF
{\True Sóng truyền từ môi trường này sang môi trường khác, đại lượng không đổi là — phát biểu đúng là: tần số.}
{Sóng truyền từ môi trường này sang môi trường khác, đại lượng không đổi là — phát biểu bước sóng là đúng.}
{\True Một dây hai đầu cố định dài $0,6\,\text{m}$ có 3 bụng. Bước sóng là — phát biểu đúng là: $0,4\,\text{m}$.}
{Một dây hai đầu cố định dài $0,6\,\text{m}$ có 3 bụng. Bước sóng là — phát biểu $1,8\,\text{m}$ là đúng.}
\loigiai{
\begin{itemchoice}
\item (Đúng) Sóng truyền từ môi trường này sang môi trường khác, đại lượng không đổi là — phát biểu đúng là: tần số. (Vì): Tần số do nguồn quyết định và không đổi khi qua mặt phân cách. b) Sai: phương án nêu không phù hợp. c) Đúng: $0,6=3\lambda/2$, suy ra $\lambda=0,4$ m. d) Sai: phương án nêu không phù hợp
\item (Sai) Sóng truyền từ môi trường này sang môi trường khác, đại lượng không đổi là — phát biểu bước sóng là đúng.
\item (Đúng) Một dây hai đầu cố định dài $0,6\,\text{m}$ có 3 bụng. Bước sóng là — phát biểu đúng là: $0,4\,\text{m}$.
\item (Sai) Một dây hai đầu cố định dài $0,6\,\text{m}$ có 3 bụng. Bước sóng là — phát biểu $1,8\,\text{m}$ là đúng.
\end{itemchoice}
}
\end{ex}

% ===== Câu 58 =====
% ID: L11C2B14-15-TL
% Mức: VD
\begin{ex}
% ID: L11C2B14-15-TL
% Mức: VD
Trong dạng “Bài toán âm, tần số và cường độ âm”, Hai nguồn cùng pha có $\lambda=3$ cm; tại $M$, $d_1=14$ cm, $d_2=21,5$ cm. Xác định trạng thái giao thoa.
\choiceTF

\loigiai{
$\Delta d=7,5\,\mathrm{cm}=2,5\lambda$, nên $M$ là cực tiểu.
}
\end{ex}

% ===== Câu 59 =====
% ID: L11C2B14-15-TLN
% Mức: VD
\begin{ex}
% ID: L11C2B14-15-TLN
% Mức: VD
Trong dạng “Bài toán âm, tần số và cường độ âm”, Hai điểm cách nhau $1\,\text{m}$ trên phương truyền sóng, $\lambda=4$ m. Tính độ lệch pha theo radian, lấy kết quả theo $\pi$.
\shortans{0,5}
\loigiai{
$\Delta\varphi=2\pi\Delta x/\lambda=2\pi/4=0,5\pi$ rad.
}
\end{ex}

% ===== Câu 60 =====
% ID: L11C2B14-15-TN
% Mức: TH
\dangbt{Bài toán sóng dừng}
\begin{ex}
% ID: L11C2B14-15-TN
% Mức: TH
Trong dạng “Bài toán sóng dừng”, Hai nút sóng dừng liên tiếp cách nhau
\choice
{\True $\lambda/2$}
{$\lambda$}
{$\lambda/4$}
{$2\lambda$}
\loigiai{
Đáp án A. Khoảng cách hai nút liên tiếp là nửa bước sóng.
}
\end{ex}

% ===== Câu 61 =====
% ID: L11C2B14-16-DS
% Mức: TH
\dangbt{Phương trình sóng và độ lệch pha}
\begin{ex}
% ID: L11C2B14-16-DS
% Mức: TH
Xét các phát biểu sau trong dạng “Phương trình sóng và độ lệch pha”.
\choiceTF
{\True Một sóng có $v=30$ m/s, $f=5$ Hz. Bước sóng là — phát biểu đúng là: $6\,\text{m}$.}
{Một sóng có $v=30$ m/s, $f=5$ Hz. Bước sóng là — phát biểu $150\,\text{m}$ là đúng.}
{\True Độ lệch pha giữa hai điểm cách nhau $\Delta x$ trên phương truyền sóng là — phát biểu đúng là: $2\pi\Delta x/\lambda$.}
{Độ lệch pha giữa hai điểm cách nhau $\Delta x$ trên phương truyền sóng là — phát biểu $2\pi\lambda/\Delta x$ là đúng.}
\loigiai{
\begin{itemchoice}
\item (Đúng) Một sóng có $v=30$ m/s, $f=5$ Hz. Bước sóng là — phát biểu đúng là: $6\,\text{m}$. (Vì): $\lambda=v/f=30/5=6$ m. b) Sai: phương án nêu không phù hợp. c) Đúng: Theo phương trình sóng, $|\Delta\varphi|=2\pi\Delta x/\lambda$. d) Sai: phương án nêu không phù hợp
\item (Sai) Một sóng có $v=30$ m/s, $f=5$ Hz. Bước sóng là — phát biểu $150\,\text{m}$ là đúng.
\item (Đúng) Độ lệch pha giữa hai điểm cách nhau $\Delta x$ trên phương truyền sóng là — phát biểu đúng là: $2\pi\Delta x/\lambda$.
\item (Sai) Độ lệch pha giữa hai điểm cách nhau $\Delta x$ trên phương truyền sóng là — phát biểu $2\pi\lambda/\Delta x$ là đúng.
\end{itemchoice}
}
\end{ex}

% ===== Câu 62 =====
% ID: L11C2B14-16-TL
% Mức: VD
\dangbt{Bài toán âm, tần số và cường độ âm}
\begin{bt}
% ID: L11C2B14-16-TL
% Mức: VD
Trong dạng “Bài toán âm, tần số và cường độ âm”, Dây hai đầu cố định dài $1,2\,\text{m}$ có tốc độ sóng $48\,\text{m}$ /s. Tính tần số họa âm thứ ba.
\loigiai{
$f_1=v/(2L)=48/2,4=20$ Hz; $f_3=3f_1=60$ Hz.
}
\end{bt}

% ===== Câu 63 =====
% ID: L11C2B14-16-TLN
% Mức: VD
\begin{ex}
% ID: L11C2B14-16-TLN
% Mức: VD
Trong dạng “Bài toán âm, tần số và cường độ âm”, Cường độ âm tăng 100 lần. Mức cường độ âm tăng bao nhiêu dB?
\shortans{20}
\loigiai{
$\Delta L=10\log100=20$ dB.
}
\end{ex}

% ===== Câu 64 =====
% ID: L11C2B14-16-TN
% Mức: VD
\dangbt{Bài toán sóng dừng}
\begin{ex}
% ID: L11C2B14-16-TN
% Mức: VD
Trong dạng “Bài toán sóng dừng”, Âm có tần số càng lớn thì cảm giác âm càng
\choice
{\True cao}
{to}
{trầm}
{nhỏ}
\loigiai{
Đáp án A. Độ cao của âm phụ thuộc chủ yếu vào tần số.
}
\end{ex}

% ===== Câu 65 =====
% ID: L11C2B14-17-DS
% Mức: TH
\dangbt{Phương trình sóng và độ lệch pha}
\begin{ex}
% ID: L11C2B14-17-DS
% Mức: TH
Xét các phát biểu sau trong dạng “Phương trình sóng và độ lệch pha”.
\choiceTF
{\True Hai điểm cách nhau $\lambda/4$ trên phương truyền sóng lệch pha — phát biểu đúng là: $\pi/2$.}
{Hai điểm cách nhau $\lambda/4$ trên phương truyền sóng lệch pha — phát biểu $\pi$ là đúng.}
{\True Hai nguồn cùng pha, cực đại giao thoa thỏa — phát biểu đúng là: $\Delta d=k\lambda$.}
{Hai nguồn cùng pha, cực đại giao thoa thỏa — phát biểu $d_1+d_2=k\lambda$ là đúng.}
\loigiai{
\begin{itemchoice}
\item (Đúng) Hai điểm cách nhau $\lambda/4$ trên phương truyền sóng lệch pha — phát biểu đúng là: $\pi/2$. (Vì): $\Delta\varphi=2\pi(\lambda/4)/\lambda=\pi/2$. b) Sai: phương án nêu không phù hợp. c) Đúng: Điều kiện cực đại của hai nguồn cùng pha là $\Delta d=k\lambda$. d) Sai: phương án nêu không phù hợp
\item (Sai) Hai điểm cách nhau $\lambda/4$ trên phương truyền sóng lệch pha — phát biểu $\pi$ là đúng.
\item (Đúng) Hai nguồn cùng pha, cực đại giao thoa thỏa — phát biểu đúng là: $\Delta d=k\lambda$.
\item (Sai) Hai nguồn cùng pha, cực đại giao thoa thỏa — phát biểu $d_1+d_2=k\lambda$ là đúng.
\end{itemchoice}
}
\end{ex}

% ===== Câu 66 =====
% ID: L11C2B14-17-TL
% Mức: VD
\dangbt{Bài toán âm, tần số và cường độ âm}
\begin{ex}
% ID: L11C2B14-17-TL
% Mức: VD
Trong dạng “Bài toán âm, tần số và cường độ âm”, Giải thích sự thay đổi bước sóng khi sóng truyền sang môi trường khác.
\choiceTF

\loigiai{
Tần số do nguồn quyết định nên không đổi; tốc độ phụ thuộc môi trường, vì $\lambda=v/f$ nên bước sóng thay đổi theo tốc độ.
}
\end{ex}

% ===== Câu 67 =====
% ID: L11C2B14-17-TLN
% Mức: VD
\begin{ex}
% ID: L11C2B14-17-TLN
% Mức: VD
Trong dạng “Bài toán âm, tần số và cường độ âm”, Dây hai đầu cố định dài $1\,\text{m}$ có bước sóng $0,5\,\text{m}$. Số bụng là bao nhiêu?
\shortans{4}
\loigiai{
$k=2L/\lambda=2/0,5=4$ bụng.
}
\end{ex}

% ===== Câu 68 =====
% ID: L11C2B14-17-TN
% Mức: VD
\dangbt{Bài toán sóng dừng}
\begin{ex}
% ID: L11C2B14-17-TN
% Mức: VD
Trong dạng “Bài toán sóng dừng”, Mức cường độ âm được tính bởi
\choice
{\True $L=10\log(I/I_0)$}
{$L=I/I_0$}
{$L=20I/I_0$}
{$L=10I_0/I$}
\loigiai{
Đáp án A. Theo định nghĩa, $L=10\log_{10}(I/I_0)$, đơn vị dB.
}
\end{ex}

% ===== Câu 69 =====
% ID: L11C2B14-18-DS
% Mức: VD
\dangbt{Phương trình sóng và độ lệch pha}
\begin{ex}
% ID: L11C2B14-18-DS
% Mức: VD
Xét các phát biểu sau trong dạng “Phương trình sóng và độ lệch pha”.
\choiceTF
{\True Hai nút sóng dừng liên tiếp cách nhau — phát biểu đúng là: $\lambda/2$.}
{Hai nút sóng dừng liên tiếp cách nhau — phát biểu $\lambda$ là đúng.}
{\True Âm có tần số càng lớn thì cảm giác âm càng — phát biểu đúng là: cao.}
{Âm có tần số càng lớn thì cảm giác âm càng — phát biểu trầm là đúng.}
\loigiai{
\begin{itemchoice}
\item (Đúng) Hai nút sóng dừng liên tiếp cách nhau — phát biểu đúng là: $\lambda/2$. (Vì): Khoảng cách hai nút liên tiếp là nửa bước sóng. b) Sai: phương án nêu không phù hợp. c) Đúng: Độ cao của âm phụ thuộc chủ yếu vào tần số. d) Sai: phương án nêu không phù hợp
\item (Sai) Hai nút sóng dừng liên tiếp cách nhau — phát biểu $\lambda$ là đúng.
\item (Đúng) Âm có tần số càng lớn thì cảm giác âm càng — phát biểu đúng là: cao.
\item (Sai) Âm có tần số càng lớn thì cảm giác âm càng — phát biểu trầm là đúng.
\end{itemchoice}
}
\end{ex}

% ===== Câu 70 =====
% ID: L11C2B14-18-TL
% Mức: VDC
\dangbt{Bài toán âm, tần số và cường độ âm}
\begin{ex}
% ID: L11C2B14-18-TL
% Mức: VDC
Trong dạng “Bài toán âm, tần số và cường độ âm”, Tại sao đứng xa nguồn âm thì âm nghe nhỏ hơn?
\choiceTF

\loigiai{
Năng lượng sóng phân bố trên diện tích ngày càng lớn; với nguồn điểm, cường độ giảm theo $1/r^2$, nên âm nghe nhỏ hơn.
}
\end{ex}

% ===== Câu 71 =====
% ID: L11C2B14-18-TLN
% Mức: VDC
\begin{ex}
% ID: L11C2B14-18-TLN
% Mức: VDC
Trong dạng “Bài toán âm, tần số và cường độ âm”, Nguồn điểm có cường độ âm 8 đơn vị tại r. Ở khoảng cách 2r, cường độ theo cùng đơn vị là bao nhiêu?
\shortans{2}
\loigiai{
$I_2=I_1(r/2r)^2=8/4=2$.
}
\end{ex}

% ===== Câu 72 =====
% ID: L11C2B14-18-TN
% Mức: VD
\dangbt{Bài toán sóng dừng}
\begin{ex}
% ID: L11C2B14-18-TN
% Mức: VD
Trong dạng “Bài toán sóng dừng”, Khi khoảng cách đến nguồn điểm tăng gấp đôi, cường độ âm trong môi trường không hấp thụ
\choice
{\True giảm 4 lần}
{giảm 2 lần}
{tăng 4 lần}
{không đổi}
\loigiai{
Đáp án A. Với nguồn điểm $I\sim1/r^2$, nên r tăng 2 lần thì $I$ giảm 4 lần.
}
\end{ex}

% ===== Câu 73 =====
% ID: L11C2B14-19-DS
% Mức: VD
\dangbt{Phương trình sóng và độ lệch pha}
\begin{ex}
% ID: L11C2B14-19-DS
% Mức: VD
Xét các phát biểu sau trong dạng “Phương trình sóng và độ lệch pha”.
\choiceTF
{\True Mức cường độ âm được tính bởi — phát biểu đúng là: $L=10\log(I/I_0)$.}
{Mức cường độ âm được tính bởi — phát biểu $L=I/I_0$ là đúng.}
{\True Khi khoảng cách đến nguồn điểm tăng gấp đôi, cường độ âm trong môi trường không hấp thụ — phát biểu đúng là: giảm 4 lần.}
{Khi khoảng cách đến nguồn điểm tăng gấp đôi, cường độ âm trong môi trường không hấp thụ — phát biểu tăng 4 lần là đúng.}
\loigiai{
\begin{itemchoice}
\item (Đúng) Mức cường độ âm được tính bởi — phát biểu đúng là: $L=10\log(I/I_0)$. (Vì): Theo định nghĩa, $L=10\log_{10}(I/I_0)$, đơn vị dB. b) Sai: phương án nêu không phù hợp. c) Đúng: Với nguồn điểm $I\sim1/r^2$, nên r tăng 2 lần thì $I$ giảm 4 lần. d) Sai: phương án nêu không phù hợp
\item (Sai) Mức cường độ âm được tính bởi — phát biểu $L=I/I_0$ là đúng.
\item (Đúng) Khi khoảng cách đến nguồn điểm tăng gấp đôi, cường độ âm trong môi trường không hấp thụ — phát biểu đúng là: giảm 4 lần.
\item (Sai) Khi khoảng cách đến nguồn điểm tăng gấp đôi, cường độ âm trong môi trường không hấp thụ — phát biểu tăng 4 lần là đúng.
\end{itemchoice}
}
\end{ex}

% ===== Câu 74 =====
% ID: L11C2B14-19-TL
% Mức: NB
\dangbt{Chưa có dạng}
\begin{ex}
% ID: L11C2B14-19-TL
% Mức: NB
Sóng dừng là A. sóng được tạo thành giữa hai điểm cố định trong một môi trường. B. sóng không lan truyền được do bị một vật cản chặn lại. C. sóng được tạo thành do sự giao thoa của hai sóng kết hợp, trên đường thẳng nối hai tâm phát sóng. $\nguon{ly11 kntt.pdf}$
\choiceTF

\loigiai{
Chọn D.

Sóng dừng là: Sóng được tạo thành do sự giao thoa giữa sóng tới và sóng phản xạ.
}
\end{ex}

% ===== Câu 75 =====
% ID: L11C2B14-19-TLN
% Mức: TH
\dangbt{Phương trình sóng và độ lệch pha}
\begin{ex}
% ID: L11C2B14-19-TLN
% Mức: TH
Trong dạng “Phương trình sóng và độ lệch pha”, Sóng có chu kì $0,25\,\text{s}$. Tính tần số theo Hz.
\shortans{4}
\loigiai{
$f=1/T=1/0,25=4$ Hz.
}
\end{ex}

% ===== Câu 76 =====
% ID: L11C2B14-19-TN
% Mức: VDC
\dangbt{Bài toán sóng dừng}
\begin{ex}
% ID: L11C2B14-19-TN
% Mức: VDC
Trong dạng “Bài toán sóng dừng”, Sóng truyền từ môi trường này sang môi trường khác, đại lượng không đổi là
\choice
{\True tần số}
{bước sóng}
{tốc độ}
{biên độ}
\loigiai{
Đáp án A. Tần số do nguồn quyết định và không đổi khi qua mặt phân cách.
}
\end{ex}

% ===== Câu 77 =====
% ID: L11C2B14-20-DS
% Mức: VDC
\dangbt{Phương trình sóng và độ lệch pha}
\begin{ex}
% ID: L11C2B14-20-DS
% Mức: VDC
Xét các phát biểu sau trong dạng “Phương trình sóng và độ lệch pha”.
\choiceTF
{\True Sóng truyền từ môi trường này sang môi trường khác, đại lượng không đổi là — phát biểu đúng là: tần số.}
{Sóng truyền từ môi trường này sang môi trường khác, đại lượng không đổi là — phát biểu bước sóng là đúng.}
{\True Một dây hai đầu cố định dài $0,6\,\text{m}$ có 3 bụng. Bước sóng là — phát biểu đúng là: $0,4\,\text{m}$.}
{Một dây hai đầu cố định dài $0,6\,\text{m}$ có 3 bụng. Bước sóng là — phát biểu $1,8\,\text{m}$ là đúng.}
\loigiai{
\begin{itemchoice}
\item (Đúng) Sóng truyền từ môi trường này sang môi trường khác, đại lượng không đổi là — phát biểu đúng là: tần số. (Vì): Tần số do nguồn quyết định và không đổi khi qua mặt phân cách. b) Sai: phương án nêu không phù hợp. c) Đúng: $0,6=3\lambda/2$, suy ra $\lambda=0,4$ m. d) Sai: phương án nêu không phù hợp
\item (Sai) Sóng truyền từ môi trường này sang môi trường khác, đại lượng không đổi là — phát biểu bước sóng là đúng.
\item (Đúng) Một dây hai đầu cố định dài $0,6\,\text{m}$ có 3 bụng. Bước sóng là — phát biểu đúng là: $0,4\,\text{m}$.
\item (Sai) Một dây hai đầu cố định dài $0,6\,\text{m}$ có 3 bụng. Bước sóng là — phát biểu $1,8\,\text{m}$ là đúng.
\end{itemchoice}
}
\end{ex}

% ===== Câu 78 =====
% ID: L11C2B14-20-TL
% Mức: NB
\dangbt{Chưa có dạng}
\begin{ex}
% ID: L11C2B14-20-TL
% Mức: NB
Một sóng ngang có tần số $20 \mathrm{Hz}$ truyền trên mặt nước với tốc độ $1,5 \mathrm{m/s}$. Trên phương truyền sóng, sóng truyền đến điểm $P$ rồi mới đến điểm $Q$ cách nó $16,125 \mathrm{cm}$. Tại thời điểm t, điểm $P$ hạ 3 xuống thấp nhất thì sau thời gian $\Delta t = 400\$, $\mathrm{s},$ điểm $Q$ sẽ tới vị trí nào? $\nguon{ly11 kntt.pdf}$
\choiceTF

\loigiai{
Tương tự Bài 9.8 ta xác định được sau thời gian $\text{\Delta t} = \dfrac{3}{400}\text{s}$ điểm $Q$ xuống vị trí thấp nhất.
}
\end{ex}

% ===== Câu 79 =====
% ID: L11C2B14-20-TLN
% Mức: TH
\dangbt{Phương trình sóng và độ lệch pha}
\begin{ex}
% ID: L11C2B14-20-TLN
% Mức: TH
Trong dạng “Phương trình sóng và độ lệch pha”, Sóng có $v=24$ m/s và $f=8$ Hz. Tính bước sóng theo m.
\shortans{3}
\loigiai{
$\lambda=v/f=24/8=3$ m.
}
\end{ex}

% ===== Câu 80 =====
% ID: L11C2B14-20-TN
% Mức: VDC
\dangbt{Bài toán sóng dừng}
\begin{ex}
% ID: L11C2B14-20-TN
% Mức: VDC
Trong dạng “Bài toán sóng dừng”, Một dây hai đầu cố định dài $0,6\,\text{m}$ có 3 bụng. Bước sóng là
\choice
{\True $0,4\,\text{m}$}
{$0,2\,\text{m}$}
{$1,8\,\text{m}$}
{$0,6\,\text{m}$}
\loigiai{
Đáp án A. $0,6=3\lambda/2$, suy ra $\lambda=0,4$ m.
}
\end{ex}

% ===== Câu 81 =====
% ID: L11C2B14-21-DS
% Mức: TH
\dangbt{Tổng hợp các đại lượng đặc trưng của sóng}
\begin{ex}
% ID: L11C2B14-21-DS
% Mức: TH
Xét các phát biểu sau trong dạng “Tổng hợp các đại lượng đặc trưng của sóng”.
\choiceTF
{\True Một sóng có $v=30$ m/s, $f=5$ Hz. Bước sóng là — phát biểu đúng là: $6\,\text{m}$.}
{Một sóng có $v=30$ m/s, $f=5$ Hz. Bước sóng là — phát biểu $150\,\text{m}$ là đúng.}
{\True Độ lệch pha giữa hai điểm cách nhau $\Delta x$ trên phương truyền sóng là — phát biểu đúng là: $2\pi\Delta x/\lambda$.}
{Độ lệch pha giữa hai điểm cách nhau $\Delta x$ trên phương truyền sóng là — phát biểu $2\pi\lambda/\Delta x$ là đúng.}
\loigiai{
\begin{itemchoice}
\item (Đúng) Một sóng có $v=30$ m/s, $f=5$ Hz. Bước sóng là — phát biểu đúng là: $6\,\text{m}$. (Vì): $\lambda=v/f=30/5=6$ m. b) Sai: phương án nêu không phù hợp. c) Đúng: Theo phương trình sóng, $|\Delta\varphi|=2\pi\Delta x/\lambda$. d) Sai: phương án nêu không phù hợp
\item (Sai) Một sóng có $v=30$ m/s, $f=5$ Hz. Bước sóng là — phát biểu $150\,\text{m}$ là đúng.
\item (Đúng) Độ lệch pha giữa hai điểm cách nhau $\Delta x$ trên phương truyền sóng là — phát biểu đúng là: $2\pi\Delta x/\lambda$.
\item (Sai) Độ lệch pha giữa hai điểm cách nhau $\Delta x$ trên phương truyền sóng là — phát biểu $2\pi\lambda/\Delta x$ là đúng.
\end{itemchoice}
}
\end{ex}

% ===== Câu 82 =====
% ID: L11C2B14-21-TL
% Mức: TH
\dangbt{Chưa có dạng}
\begin{ex}
% ID: L11C2B14-21-TL
% Mức: TH
Một sợi dây đàn hồi, mảnh, dài có đầu $O$ dao động với tần số f thay đổi được trong khoảng từ $80 \mathrm{Hz}$ đến $125 \mathrm{Hz}$, theo phương vuông góc với sợi dây. Sóng tạo thành lan truyền trên dây với tốc độ không đổi $v =$ 10 $\mathrm{m/s}$. a) Cho $f =$ 80 $\mathrm{Hz},$ tính chu kì và bước sóng của sóng trên dây. b) Tính tần số f để điểm $M$ trên dây cách $O$ một khoảng bằng $20 \mathrm{cm}$ luôn dao động cùng pha với điểm O. $\nguon{ly11 kntt.pdf}$
\choiceTF

\loigiai{
a) $T = \dfrac{1}{f} = \dfrac{1}{80} = 0,0125\text{s};\lambda = \dfrac{v}{f} = \dfrac{10}{80} = 0,125\text{m}$.

b) Vì sóng tại hai điểm $M$, $O$ cùng pha nhau, nên khoảng cách OM thoả mãn:

OM = d = kλ = k $\dfrac{v}{f}\Rightarrow$  f = $\dfrac{kv}{d} = \dfrac{k \cdot 10}{0,2} = 50k$, với k∈Z.

Theo đề bài: $80\,\mathrm{Hz}$ ≤f≤ $125\,\mathrm{Hz}\Rightarrow80≤50k≤125\Rightarrow 1,6≤k≤2,5

Vậy k = 2. Suy ra tần số sóng là: f = 50.2 =$ 100\,\mathrm{Hz}.
}
\end{ex}

% ===== Câu 83 =====
% ID: L11C2B14-21-TLN
% Mức: VD
\dangbt{Phương trình sóng và độ lệch pha}
\begin{ex}
% ID: L11C2B14-21-TLN
% Mức: VD
Trong dạng “Phương trình sóng và độ lệch pha”, Hai điểm cách nhau $1\,\text{m}$ trên phương truyền sóng, $\lambda=4$ m. Tính độ lệch pha theo radian, lấy kết quả theo $\pi$.
\shortans{0,5}
\loigiai{
$\Delta\varphi=2\pi\Delta x/\lambda=2\pi/4=0,5\pi$ rad.
}
\end{ex}

% ===== Câu 84 =====
% ID: L11C2B14-21-TN
% Mức: NB
\dangbt{Bài toán âm, tần số và cường độ âm}
\begin{ex}
% ID: L11C2B14-21-TN
% Mức: NB
Trong dạng “Bài toán âm, tần số và cường độ âm”, Một sóng có $v=30$ m/s, $f=5$ Hz. Bước sóng là
\choice
{\True $6\,\text{m}$}
{$150\,\text{m}$}
{$0,167\,\text{m}$}
{$35\,\text{m}$}
\loigiai{
Đáp án A. $\lambda=v/f=30/5=6$ m.
}
\end{ex}

% ===== Câu 85 =====
% ID: L11C2B14-22-DS
% Mức: TH
\dangbt{Tổng hợp các đại lượng đặc trưng của sóng}
\begin{ex}
% ID: L11C2B14-22-DS
% Mức: TH
Xét các phát biểu sau trong dạng “Tổng hợp các đại lượng đặc trưng của sóng”.
\choiceTF
{\True Hai điểm cách nhau $\lambda/4$ trên phương truyền sóng lệch pha — phát biểu đúng là: $\pi/2$.}
{Hai điểm cách nhau $\lambda/4$ trên phương truyền sóng lệch pha — phát biểu $\pi$ là đúng.}
{\True Hai nguồn cùng pha, cực đại giao thoa thỏa — phát biểu đúng là: $\Delta d=k\lambda$.}
{Hai nguồn cùng pha, cực đại giao thoa thỏa — phát biểu $d_1+d_2=k\lambda$ là đúng.}
\loigiai{
\begin{itemchoice}
\item (Đúng) Hai điểm cách nhau $\lambda/4$ trên phương truyền sóng lệch pha — phát biểu đúng là: $\pi/2$. (Vì): $\Delta\varphi=2\pi(\lambda/4)/\lambda=\pi/2$. b) Sai: phương án nêu không phù hợp. c) Đúng: Điều kiện cực đại của hai nguồn cùng pha là $\Delta d=k\lambda$. d) Sai: phương án nêu không phù hợp
\item (Sai) Hai điểm cách nhau $\lambda/4$ trên phương truyền sóng lệch pha — phát biểu $\pi$ là đúng.
\item (Đúng) Hai nguồn cùng pha, cực đại giao thoa thỏa — phát biểu đúng là: $\Delta d=k\lambda$.
\item (Sai) Hai nguồn cùng pha, cực đại giao thoa thỏa — phát biểu $d_1+d_2=k\lambda$ là đúng.
\end{itemchoice}
}
\end{ex}

% ===== Câu 86 =====
% ID: L11C2B14-22-TL
% Mức: TH
\dangbt{Chưa có dạng}
\begin{ex}
% ID: L11C2B14-22-TL
% Mức: TH
Sóng vô tuyến ngắn có thể được sử dụng để đo khoảng cách từ Trái Đất đến Mặt Trăng, bằng cách phát một tín hiệu từ Trái Đất tới Mặt Trăng và thu tín hiệu trở lại, đo khoảng thời gian từ khi phát đến khi nhận tín hiệu. a) Biết khoảng thời gian từ khi phát tới khi nhận được tín hiệu trở lại là $2,6\,\mathrm{s}$. Tính khoảng cách từ Mặt Trăng tới Trái Đất. b) Sóng vô tuyến trên có tần số $107 \mathrm{Hz}$. Tính bước sóng của sóng. $\nguon{ly11 kntt.pdf}$
\choiceTF

\loigiai{
Chọn D.

Sóng dừng là: Sóng được tạo thành do sự giao thoa giữa sóng tới và sóng phản xạ.
}
\end{ex}

% ===== Câu 87 =====
% ID: L11C2B14-22-TLN
% Mức: VD
\dangbt{Phương trình sóng và độ lệch pha}
\begin{ex}
% ID: L11C2B14-22-TLN
% Mức: VD
Trong dạng “Phương trình sóng và độ lệch pha”, Cường độ âm tăng 100 lần. Mức cường độ âm tăng bao nhiêu dB?
\shortans{20}
\loigiai{
$\Delta L=10\log100=20$ dB.
}
\end{ex}

% ===== Câu 88 =====
% ID: L11C2B14-22-TN
% Mức: NB
\dangbt{Bài toán âm, tần số và cường độ âm}
\begin{ex}
% ID: L11C2B14-22-TN
% Mức: NB
Trong dạng “Bài toán âm, tần số và cường độ âm”, Độ lệch pha giữa hai điểm cách nhau $\Delta x$ trên phương truyền sóng là
\choice
{\True $2\pi\Delta x/\lambda$}
{$\pi\Delta x/\lambda$}
{$2\pi\lambda/\Delta x$}
{$\Delta x/(2\pi\lambda)$}
\loigiai{
Đáp án A. Theo phương trình sóng, $|\Delta\varphi|=2\pi\Delta x/\lambda$.
}
\end{ex}

% ===== Câu 89 =====
% ID: L11C2B14-23-DS
% Mức: VD
\dangbt{Tổng hợp các đại lượng đặc trưng của sóng}
\begin{ex}
% ID: L11C2B14-23-DS
% Mức: VD
Xét các phát biểu sau trong dạng “Tổng hợp các đại lượng đặc trưng của sóng”.
\choiceTF
{\True Hai nút sóng dừng liên tiếp cách nhau — phát biểu đúng là: $\lambda/2$.}
{Hai nút sóng dừng liên tiếp cách nhau — phát biểu $\lambda$ là đúng.}
{\True Âm có tần số càng lớn thì cảm giác âm càng — phát biểu đúng là: cao.}
{Âm có tần số càng lớn thì cảm giác âm càng — phát biểu trầm là đúng.}
\loigiai{
\begin{itemchoice}
\item (Đúng) Hai nút sóng dừng liên tiếp cách nhau — phát biểu đúng là: $\lambda/2$. (Vì): Khoảng cách hai nút liên tiếp là nửa bước sóng. b) Sai: phương án nêu không phù hợp. c) Đúng: Độ cao của âm phụ thuộc chủ yếu vào tần số. d) Sai: phương án nêu không phù hợp
\item (Sai) Hai nút sóng dừng liên tiếp cách nhau — phát biểu $\lambda$ là đúng.
\item (Đúng) Âm có tần số càng lớn thì cảm giác âm càng — phát biểu đúng là: cao.
\item (Sai) Âm có tần số càng lớn thì cảm giác âm càng — phát biểu trầm là đúng.
\end{itemchoice}
}
\end{ex}

% ===== Câu 90 =====
% ID: L11C2B14-23-TL
% Mức: TH
\dangbt{Chưa có dạng}
\begin{ex}
% ID: L11C2B14-23-TL
% Mức: TH
Trong thí nghiệm Young về giao thoa với ánh sáng đơn sắc. Trên màn chỉ quan sát được 21 vạch sáng mà khoảng cách giữa hai vạch sáng ngoài cùng là $4 \mathrm{cm}$. Tại hai điểm $P$ và $Q$ là hai vị trí cho vân sáng trên màn. Hãy xác định số vân sáng trên đoạn PQ, biết rằng khoảng cách giữa hai điểm đó là $2,4 \mathrm{cm}$. CHƯƠNG III. ĐIỆ $N$ TRƯỜNG BÀ $I$ 16. LỰ $C$ TƯONG TÁ $C$ GIŨ $A$ CÁ $C$ ĐIỆ $N$ TÍCH $\nguon{ly11 kntt.pdf}$
\choiceTF

\loigiai{
$\left. i = \dfrac{\text{\Delta }s}{21 - 1} = \dfrac{40}{20} = 2\text{mm}\Rightarrow N_{s} = \left( \dfrac{PQ}{i} \right) + 1 = 12 + 1 = 13 \right.$.
}
\end{ex}

% ===== Câu 91 =====
% ID: L11C2B14-23-TLN
% Mức: VD
\dangbt{Phương trình sóng và độ lệch pha}
\begin{ex}
% ID: L11C2B14-23-TLN
% Mức: VD
Trong dạng “Phương trình sóng và độ lệch pha”, Dây hai đầu cố định dài $1\,\text{m}$ có bước sóng $0,5\,\text{m}$. Số bụng là bao nhiêu?
\shortans{4}
\loigiai{
$k=2L/\lambda=2/0,5=4$ bụng.
}
\end{ex}

% ===== Câu 92 =====
% ID: L11C2B14-23-TN
% Mức: TH
\dangbt{Bài toán âm, tần số và cường độ âm}
\begin{ex}
% ID: L11C2B14-23-TN
% Mức: TH
Trong dạng “Bài toán âm, tần số và cường độ âm”, Hai điểm cách nhau $\lambda/4$ trên phương truyền sóng lệch pha
\choice
{\True $\pi/2$}
{$\pi$}
{$2\pi$}
{$\pi/4$}
\loigiai{
Đáp án A. $\Delta\varphi=2\pi(\lambda/4)/\lambda=\pi/2$.
}
\end{ex}

% ===== Câu 93 =====
% ID: L11C2B14-24-DS
% Mức: VD
\dangbt{Tổng hợp các đại lượng đặc trưng của sóng}
\begin{ex}
% ID: L11C2B14-24-DS
% Mức: VD
Xét các phát biểu sau trong dạng “Tổng hợp các đại lượng đặc trưng của sóng”.
\choiceTF
{\True Mức cường độ âm được tính bởi — phát biểu đúng là: $L=10\log(I/I_0)$.}
{Mức cường độ âm được tính bởi — phát biểu $L=I/I_0$ là đúng.}
{\True Khi khoảng cách đến nguồn điểm tăng gấp đôi, cường độ âm trong môi trường không hấp thụ — phát biểu đúng là: giảm 4 lần.}
{Khi khoảng cách đến nguồn điểm tăng gấp đôi, cường độ âm trong môi trường không hấp thụ — phát biểu tăng 4 lần là đúng.}
\loigiai{
\begin{itemchoice}
\item (Đúng) Mức cường độ âm được tính bởi — phát biểu đúng là: $L=10\log(I/I_0)$. (Vì): Theo định nghĩa, $L=10\log_{10}(I/I_0)$, đơn vị dB. b) Sai: phương án nêu không phù hợp. c) Đúng: Với nguồn điểm $I\sim1/r^2$, nên r tăng 2 lần thì $I$ giảm 4 lần. d) Sai: phương án nêu không phù hợp
\item (Sai) Mức cường độ âm được tính bởi — phát biểu $L=I/I_0$ là đúng.
\item (Đúng) Khi khoảng cách đến nguồn điểm tăng gấp đôi, cường độ âm trong môi trường không hấp thụ — phát biểu đúng là: giảm 4 lần.
\item (Sai) Khi khoảng cách đến nguồn điểm tăng gấp đôi, cường độ âm trong môi trường không hấp thụ — phát biểu tăng 4 lần là đúng.
\end{itemchoice}
}
\end{ex}

% ===== Câu 94 =====
% ID: L11C2B14-24-TL
% Mức: TH
\dangbt{Phương trình sóng và độ lệch pha}
\begin{bt}
% ID: L11C2B14-24-TL
% Mức: TH
Trong dạng “Phương trình sóng và độ lệch pha”, Một sóng truyền với tốc độ $20\,\text{m}$ /s và chu kì $0,1\,\text{s}$. Tính bước sóng.
\loigiai{
$f=1/T=10$ Hz; $\lambda=v/f=20/10=2$ m.
}
\end{bt}

% ===== Câu 95 =====
% ID: L11C2B14-24-TLN
% Mức: VDC
\begin{ex}
% ID: L11C2B14-24-TLN
% Mức: VDC
Trong dạng “Phương trình sóng và độ lệch pha”, Nguồn điểm có cường độ âm 8 đơn vị tại r. Ở khoảng cách 2r, cường độ theo cùng đơn vị là bao nhiêu?
\shortans{2}
\loigiai{
$I_2=I_1(r/2r)^2=8/4=2$.
}
\end{ex}

% ===== Câu 96 =====
% ID: L11C2B14-24-TN
% Mức: TH
\dangbt{Bài toán âm, tần số và cường độ âm}
\begin{ex}
% ID: L11C2B14-24-TN
% Mức: TH
Trong dạng “Bài toán âm, tần số và cường độ âm”, Hai nguồn cùng pha, cực đại giao thoa thỏa
\choice
{\True $\Delta d=k\lambda$}
{$\Delta d=(k+1/2)\lambda$}
{$d_1+d_2=k\lambda$}
{$\Delta d=\lambda/4$}
\loigiai{
Đáp án A. Điều kiện cực đại của hai nguồn cùng pha là $\Delta d=k\lambda$.
}
\end{ex}

% ===== Câu 97 =====
% ID: L11C2B14-25-DS
% Mức: VDC
\dangbt{Tổng hợp các đại lượng đặc trưng của sóng}
\begin{ex}
% ID: L11C2B14-25-DS
% Mức: VDC
Xét các phát biểu sau trong dạng “Tổng hợp các đại lượng đặc trưng của sóng”.
\choiceTF
{\True Sóng truyền từ môi trường này sang môi trường khác, đại lượng không đổi là — phát biểu đúng là: tần số.}
{Sóng truyền từ môi trường này sang môi trường khác, đại lượng không đổi là — phát biểu bước sóng là đúng.}
{\True Một dây hai đầu cố định dài $0,6\,\text{m}$ có 3 bụng. Bước sóng là — phát biểu đúng là: $0,4\,\text{m}$.}
{Một dây hai đầu cố định dài $0,6\,\text{m}$ có 3 bụng. Bước sóng là — phát biểu $1,8\,\text{m}$ là đúng.}
\loigiai{
\begin{itemchoice}
\item (Đúng) Sóng truyền từ môi trường này sang môi trường khác, đại lượng không đổi là — phát biểu đúng là: tần số. (Vì): Tần số do nguồn quyết định và không đổi khi qua mặt phân cách. b) Sai: phương án nêu không phù hợp. c) Đúng: $0,6=3\lambda/2$, suy ra $\lambda=0,4$ m. d) Sai: phương án nêu không phù hợp
\item (Sai) Sóng truyền từ môi trường này sang môi trường khác, đại lượng không đổi là — phát biểu bước sóng là đúng.
\item (Đúng) Một dây hai đầu cố định dài $0,6\,\text{m}$ có 3 bụng. Bước sóng là — phát biểu đúng là: $0,4\,\text{m}$.
\item (Sai) Một dây hai đầu cố định dài $0,6\,\text{m}$ có 3 bụng. Bước sóng là — phát biểu $1,8\,\text{m}$ là đúng.
\end{itemchoice}
}
\end{ex}

% ===== Câu 98 =====
% ID: L11C2B14-25-TL
% Mức: TH
\dangbt{Phương trình sóng và độ lệch pha}
\begin{ex}
% ID: L11C2B14-25-TL
% Mức: TH
Trong dạng “Phương trình sóng và độ lệch pha”, Viết biểu thức độ lệch pha của hai điểm trên phương truyền sóng và nêu ý nghĩa.
\choiceTF

\loigiai{
$\Delta\varphi=2\pi\Delta x/\lambda$. Công thức cho biết quan hệ pha phụ thuộc khoảng cách theo phương truyền sóng.
}
\end{ex}

% ===== Câu 99 =====
% ID: L11C2B14-25-TLN
% Mức: TH
\dangbt{Tổng hợp các đại lượng đặc trưng của sóng}
\begin{ex}
% ID: L11C2B14-25-TLN
% Mức: TH
Trong dạng “Tổng hợp các đại lượng đặc trưng của sóng”, Sóng có chu kì $0,25\,\text{s}$. Tính tần số theo Hz.
\shortans{4}
\loigiai{
$f=1/T=1/0,25=4$ Hz.
}
\end{ex}

% ===== Câu 100 =====
% ID: L11C2B14-25-TN
% Mức: TH
\dangbt{Bài toán âm, tần số và cường độ âm}
\begin{ex}
% ID: L11C2B14-25-TN
% Mức: TH
Trong dạng “Bài toán âm, tần số và cường độ âm”, Hai nút sóng dừng liên tiếp cách nhau
\choice
{\True $\lambda/2$}
{$\lambda$}
{$\lambda/4$}
{$2\lambda$}
\loigiai{
Đáp án A. Khoảng cách hai nút liên tiếp là nửa bước sóng.
}
\end{ex}

% ===== Câu 101 =====
% ID: L11C2B14-26-DS
% Mức: TH
\dangbt{Vận dụng tổng hợp kiến thức sóng trong thực tế}
\begin{ex}
% ID: L11C2B14-26-DS
% Mức: TH
Xét các phát biểu sau trong dạng “Vận dụng tổng hợp kiến thức sóng trong thực tế”.
\choiceTF
{\True Một sóng có $v=30$ m/s, $f=5$ Hz. Bước sóng là — phát biểu đúng là: $6\,\text{m}$.}
{Một sóng có $v=30$ m/s, $f=5$ Hz. Bước sóng là — phát biểu $150\,\text{m}$ là đúng.}
{\True Độ lệch pha giữa hai điểm cách nhau $\Delta x$ trên phương truyền sóng là — phát biểu đúng là: $2\pi\Delta x/\lambda$.}
{Độ lệch pha giữa hai điểm cách nhau $\Delta x$ trên phương truyền sóng là — phát biểu $2\pi\lambda/\Delta x$ là đúng.}
\loigiai{
\begin{itemchoice}
\item (Đúng) Một sóng có $v=30$ m/s, $f=5$ Hz. Bước sóng là — phát biểu đúng là: $6\,\text{m}$. (Vì): $\lambda=v/f=30/5=6$ m. b) Sai: phương án nêu không phù hợp. c) Đúng: Theo phương trình sóng, $|\Delta\varphi|=2\pi\Delta x/\lambda$. d) Sai: phương án nêu không phù hợp
\item (Sai) Một sóng có $v=30$ m/s, $f=5$ Hz. Bước sóng là — phát biểu $150\,\text{m}$ là đúng.
\item (Đúng) Độ lệch pha giữa hai điểm cách nhau $\Delta x$ trên phương truyền sóng là — phát biểu đúng là: $2\pi\Delta x/\lambda$.
\item (Sai) Độ lệch pha giữa hai điểm cách nhau $\Delta x$ trên phương truyền sóng là — phát biểu $2\pi\lambda/\Delta x$ là đúng.
\end{itemchoice}
}
\end{ex}

% ===== Câu 102 =====
% ID: L11C2B14-26-TL
% Mức: VD
\dangbt{Phương trình sóng và độ lệch pha}
\begin{ex}
% ID: L11C2B14-26-TL
% Mức: VD
Trong dạng “Phương trình sóng và độ lệch pha”, Hai nguồn cùng pha có $\lambda=3$ cm; tại $M$, $d_1=14$ cm, $d_2=21,5$ cm. Xác định trạng thái giao thoa.
\choiceTF

\loigiai{
$\Delta d=7,5\,\mathrm{cm}=2,5\lambda$, nên $M$ là cực tiểu.
}
\end{ex}

% ===== Câu 103 =====
% ID: L11C2B14-26-TLN
% Mức: TH
\dangbt{Tổng hợp các đại lượng đặc trưng của sóng}
\begin{ex}
% ID: L11C2B14-26-TLN
% Mức: TH
Trong dạng “Tổng hợp các đại lượng đặc trưng của sóng”, Sóng có $v=24$ m/s và $f=8$ Hz. Tính bước sóng theo m.
\shortans{3}
\loigiai{
$\lambda=v/f=24/8=3$ m.
}
\end{ex}

% ===== Câu 104 =====
% ID: L11C2B14-26-TN
% Mức: VD
\dangbt{Bài toán âm, tần số và cường độ âm}
\begin{ex}
% ID: L11C2B14-26-TN
% Mức: VD
Trong dạng “Bài toán âm, tần số và cường độ âm”, Âm có tần số càng lớn thì cảm giác âm càng
\choice
{\True cao}
{to}
{trầm}
{nhỏ}
\loigiai{
Đáp án A. Độ cao của âm phụ thuộc chủ yếu vào tần số.
}
\end{ex}

% ===== Câu 105 =====
% ID: L11C2B14-27-DS
% Mức: TH
\dangbt{Vận dụng tổng hợp kiến thức sóng trong thực tế}
\begin{ex}
% ID: L11C2B14-27-DS
% Mức: TH
Xét các phát biểu sau trong dạng “Vận dụng tổng hợp kiến thức sóng trong thực tế”.
\choiceTF
{\True Hai điểm cách nhau $\lambda/4$ trên phương truyền sóng lệch pha — phát biểu đúng là: $\pi/2$.}
{Hai điểm cách nhau $\lambda/4$ trên phương truyền sóng lệch pha — phát biểu $\pi$ là đúng.}
{\True Hai nguồn cùng pha, cực đại giao thoa thỏa — phát biểu đúng là: $\Delta d=k\lambda$.}
{Hai nguồn cùng pha, cực đại giao thoa thỏa — phát biểu $d_1+d_2=k\lambda$ là đúng.}
\loigiai{
\begin{itemchoice}
\item (Đúng) Hai điểm cách nhau $\lambda/4$ trên phương truyền sóng lệch pha — phát biểu đúng là: $\pi/2$. (Vì): $\Delta\varphi=2\pi(\lambda/4)/\lambda=\pi/2$. b) Sai: phương án nêu không phù hợp. c) Đúng: Điều kiện cực đại của hai nguồn cùng pha là $\Delta d=k\lambda$. d) Sai: phương án nêu không phù hợp
\item (Sai) Hai điểm cách nhau $\lambda/4$ trên phương truyền sóng lệch pha — phát biểu $\pi$ là đúng.
\item (Đúng) Hai nguồn cùng pha, cực đại giao thoa thỏa — phát biểu đúng là: $\Delta d=k\lambda$.
\item (Sai) Hai nguồn cùng pha, cực đại giao thoa thỏa — phát biểu $d_1+d_2=k\lambda$ là đúng.
\end{itemchoice}
}
\end{ex}

% ===== Câu 106 =====
% ID: L11C2B14-27-TL
% Mức: VD
\dangbt{Phương trình sóng và độ lệch pha}
\begin{bt}
% ID: L11C2B14-27-TL
% Mức: VD
Trong dạng “Phương trình sóng và độ lệch pha”, Dây hai đầu cố định dài $1,2\,\text{m}$ có tốc độ sóng $48\,\text{m}$ /s. Tính tần số họa âm thứ ba.
\loigiai{
$f_1=v/(2L)=48/2,4=20$ Hz; $f_3=3f_1=60$ Hz.
}
\end{bt}

% ===== Câu 107 =====
% ID: L11C2B14-27-TLN
% Mức: VD
\dangbt{Tổng hợp các đại lượng đặc trưng của sóng}
\begin{ex}
% ID: L11C2B14-27-TLN
% Mức: VD
Trong dạng “Tổng hợp các đại lượng đặc trưng của sóng”, Hai điểm cách nhau $1\,\text{m}$ trên phương truyền sóng, $\lambda=4$ m. Tính độ lệch pha theo radian, lấy kết quả theo $\pi$.
\shortans{0,5}
\loigiai{
$\Delta\varphi=2\pi\Delta x/\lambda=2\pi/4=0,5\pi$ rad.
}
\end{ex}

% ===== Câu 108 =====
% ID: L11C2B14-27-TN
% Mức: VD
\dangbt{Bài toán âm, tần số và cường độ âm}
\begin{ex}
% ID: L11C2B14-27-TN
% Mức: VD
Trong dạng “Bài toán âm, tần số và cường độ âm”, Mức cường độ âm được tính bởi
\choice
{\True $L=10\log(I/I_0)$}
{$L=I/I_0$}
{$L=20I/I_0$}
{$L=10I_0/I$}
\loigiai{
Đáp án A. Theo định nghĩa, $L=10\log_{10}(I/I_0)$, đơn vị dB.
}
\end{ex}

% ===== Câu 109 =====
% ID: L11C2B14-28-DS
% Mức: VD
\dangbt{Vận dụng tổng hợp kiến thức sóng trong thực tế}
\begin{ex}
% ID: L11C2B14-28-DS
% Mức: VD
Xét các phát biểu sau trong dạng “Vận dụng tổng hợp kiến thức sóng trong thực tế”.
\choiceTF
{\True Hai nút sóng dừng liên tiếp cách nhau — phát biểu đúng là: $\lambda/2$.}
{Hai nút sóng dừng liên tiếp cách nhau — phát biểu $\lambda$ là đúng.}
{\True Âm có tần số càng lớn thì cảm giác âm càng — phát biểu đúng là: cao.}
{Âm có tần số càng lớn thì cảm giác âm càng — phát biểu trầm là đúng.}
\loigiai{
\begin{itemchoice}
\item (Đúng) Hai nút sóng dừng liên tiếp cách nhau — phát biểu đúng là: $\lambda/2$. (Vì): Khoảng cách hai nút liên tiếp là nửa bước sóng. b) Sai: phương án nêu không phù hợp. c) Đúng: Độ cao của âm phụ thuộc chủ yếu vào tần số. d) Sai: phương án nêu không phù hợp
\item (Sai) Hai nút sóng dừng liên tiếp cách nhau — phát biểu $\lambda$ là đúng.
\item (Đúng) Âm có tần số càng lớn thì cảm giác âm càng — phát biểu đúng là: cao.
\item (Sai) Âm có tần số càng lớn thì cảm giác âm càng — phát biểu trầm là đúng.
\end{itemchoice}
}
\end{ex}

% ===== Câu 110 =====
% ID: L11C2B14-28-TL
% Mức: VD
\dangbt{Phương trình sóng và độ lệch pha}
\begin{ex}
% ID: L11C2B14-28-TL
% Mức: VD
Trong dạng “Phương trình sóng và độ lệch pha”, Giải thích sự thay đổi bước sóng khi sóng truyền sang môi trường khác.
\choiceTF

\loigiai{
Tần số do nguồn quyết định nên không đổi; tốc độ phụ thuộc môi trường, vì $\lambda=v/f$ nên bước sóng thay đổi theo tốc độ.
}
\end{ex}

% ===== Câu 111 =====
% ID: L11C2B14-28-TLN
% Mức: VD
\dangbt{Tổng hợp các đại lượng đặc trưng của sóng}
\begin{ex}
% ID: L11C2B14-28-TLN
% Mức: VD
Trong dạng “Tổng hợp các đại lượng đặc trưng của sóng”, Cường độ âm tăng 100 lần. Mức cường độ âm tăng bao nhiêu dB?
\shortans{20}
\loigiai{
$\Delta L=10\log100=20$ dB.
}
\end{ex}

% ===== Câu 112 =====
% ID: L11C2B14-28-TN
% Mức: VD
\dangbt{Bài toán âm, tần số và cường độ âm}
\begin{ex}
% ID: L11C2B14-28-TN
% Mức: VD
Trong dạng “Bài toán âm, tần số và cường độ âm”, Khi khoảng cách đến nguồn điểm tăng gấp đôi, cường độ âm trong môi trường không hấp thụ
\choice
{\True giảm 4 lần}
{giảm 2 lần}
{tăng 4 lần}
{không đổi}
\loigiai{
Đáp án A. Với nguồn điểm $I\sim1/r^2$, nên r tăng 2 lần thì $I$ giảm 4 lần.
}
\end{ex}

% ===== Câu 113 =====
% ID: L11C2B14-29-DS
% Mức: VD
\dangbt{Vận dụng tổng hợp kiến thức sóng trong thực tế}
\begin{ex}
% ID: L11C2B14-29-DS
% Mức: VD
Xét các phát biểu sau trong dạng “Vận dụng tổng hợp kiến thức sóng trong thực tế”.
\choiceTF
{\True Mức cường độ âm được tính bởi — phát biểu đúng là: $L=10\log(I/I_0)$.}
{Mức cường độ âm được tính bởi — phát biểu $L=I/I_0$ là đúng.}
{\True Khi khoảng cách đến nguồn điểm tăng gấp đôi, cường độ âm trong môi trường không hấp thụ — phát biểu đúng là: giảm 4 lần.}
{Khi khoảng cách đến nguồn điểm tăng gấp đôi, cường độ âm trong môi trường không hấp thụ — phát biểu tăng 4 lần là đúng.}
\loigiai{
\begin{itemchoice}
\item (Đúng) Mức cường độ âm được tính bởi — phát biểu đúng là: $L=10\log(I/I_0)$. (Vì): Theo định nghĩa, $L=10\log_{10}(I/I_0)$, đơn vị dB. b) Sai: phương án nêu không phù hợp. c) Đúng: Với nguồn điểm $I\sim1/r^2$, nên r tăng 2 lần thì $I$ giảm 4 lần. d) Sai: phương án nêu không phù hợp
\item (Sai) Mức cường độ âm được tính bởi — phát biểu $L=I/I_0$ là đúng.
\item (Đúng) Khi khoảng cách đến nguồn điểm tăng gấp đôi, cường độ âm trong môi trường không hấp thụ — phát biểu đúng là: giảm 4 lần.
\item (Sai) Khi khoảng cách đến nguồn điểm tăng gấp đôi, cường độ âm trong môi trường không hấp thụ — phát biểu tăng 4 lần là đúng.
\end{itemchoice}
}
\end{ex}

% ===== Câu 114 =====
% ID: L11C2B14-29-TL
% Mức: VDC
\dangbt{Phương trình sóng và độ lệch pha}
\begin{ex}
% ID: L11C2B14-29-TL
% Mức: VDC
Trong dạng “Phương trình sóng và độ lệch pha”, Tại sao đứng xa nguồn âm thì âm nghe nhỏ hơn?
\choiceTF

\loigiai{
Năng lượng sóng phân bố trên diện tích ngày càng lớn; với nguồn điểm, cường độ giảm theo $1/r^2$, nên âm nghe nhỏ hơn.
}
\end{ex}

% ===== Câu 115 =====
% ID: L11C2B14-29-TLN
% Mức: VD
\dangbt{Tổng hợp các đại lượng đặc trưng của sóng}
\begin{ex}
% ID: L11C2B14-29-TLN
% Mức: VD
Trong dạng “Tổng hợp các đại lượng đặc trưng của sóng”, Dây hai đầu cố định dài $1\,\text{m}$ có bước sóng $0,5\,\text{m}$. Số bụng là bao nhiêu?
\shortans{4}
\loigiai{
$k=2L/\lambda=2/0,5=4$ bụng.
}
\end{ex}

% ===== Câu 116 =====
% ID: L11C2B14-29-TN
% Mức: VDC
\dangbt{Bài toán âm, tần số và cường độ âm}
\begin{ex}
% ID: L11C2B14-29-TN
% Mức: VDC
Trong dạng “Bài toán âm, tần số và cường độ âm”, Sóng truyền từ môi trường này sang môi trường khác, đại lượng không đổi là
\choice
{\True tần số}
{bước sóng}
{tốc độ}
{biên độ}
\loigiai{
Đáp án A. Tần số do nguồn quyết định và không đổi khi qua mặt phân cách.
}
\end{ex}

% ===== Câu 117 =====
% ID: L11C2B14-30-DS
% Mức: VDC
\dangbt{Vận dụng tổng hợp kiến thức sóng trong thực tế}
\begin{ex}
% ID: L11C2B14-30-DS
% Mức: VDC
Xét các phát biểu sau trong dạng “Vận dụng tổng hợp kiến thức sóng trong thực tế”.
\choiceTF
{\True Sóng truyền từ môi trường này sang môi trường khác, đại lượng không đổi là — phát biểu đúng là: tần số.}
{Sóng truyền từ môi trường này sang môi trường khác, đại lượng không đổi là — phát biểu bước sóng là đúng.}
{\True Một dây hai đầu cố định dài $0,6\,\text{m}$ có 3 bụng. Bước sóng là — phát biểu đúng là: $0,4\,\text{m}$.}
{Một dây hai đầu cố định dài $0,6\,\text{m}$ có 3 bụng. Bước sóng là — phát biểu $1,8\,\text{m}$ là đúng.}
\loigiai{
\begin{itemchoice}
\item (Đúng) Sóng truyền từ môi trường này sang môi trường khác, đại lượng không đổi là — phát biểu đúng là: tần số. (Vì): Tần số do nguồn quyết định và không đổi khi qua mặt phân cách. b) Sai: phương án nêu không phù hợp. c) Đúng: $0,6=3\lambda/2$, suy ra $\lambda=0,4$ m. d) Sai: phương án nêu không phù hợp
\item (Sai) Sóng truyền từ môi trường này sang môi trường khác, đại lượng không đổi là — phát biểu bước sóng là đúng.
\item (Đúng) Một dây hai đầu cố định dài $0,6\,\text{m}$ có 3 bụng. Bước sóng là — phát biểu đúng là: $0,4\,\text{m}$.
\item (Sai) Một dây hai đầu cố định dài $0,6\,\text{m}$ có 3 bụng. Bước sóng là — phát biểu $1,8\,\text{m}$ là đúng.
\end{itemchoice}
}
\end{ex}

% ===== Câu 118 =====
% ID: L11C2B14-30-TL
% Mức: TH
\dangbt{Tổng hợp các đại lượng đặc trưng của sóng}
\begin{bt}
% ID: L11C2B14-30-TL
% Mức: TH
Trong dạng “Tổng hợp các đại lượng đặc trưng của sóng”, Một sóng truyền với tốc độ $20\,\text{m}$ /s và chu kì $0,1\,\text{s}$. Tính bước sóng.
\loigiai{
$f=1/T=10$ Hz; $\lambda=v/f=20/10=2$ m.
}
\end{bt}

% ===== Câu 119 =====
% ID: L11C2B14-30-TLN
% Mức: VDC
\begin{ex}
% ID: L11C2B14-30-TLN
% Mức: VDC
Trong dạng “Tổng hợp các đại lượng đặc trưng của sóng”, Nguồn điểm có cường độ âm 8 đơn vị tại r. Ở khoảng cách 2r, cường độ theo cùng đơn vị là bao nhiêu?
\shortans{2}
\loigiai{
$I_2=I_1(r/2r)^2=8/4=2$.
}
\end{ex}

% ===== Câu 120 =====
% ID: L11C2B14-30-TN
% Mức: VDC
\dangbt{Bài toán âm, tần số và cường độ âm}
\begin{ex}
% ID: L11C2B14-30-TN
% Mức: VDC
Trong dạng “Bài toán âm, tần số và cường độ âm”, Một dây hai đầu cố định dài $0,6\,\text{m}$ có 3 bụng. Bước sóng là
\choice
{\True $0,4\,\text{m}$}
{$0,2\,\text{m}$}
{$1,8\,\text{m}$}
{$0,6\,\text{m}$}
\loigiai{
Đáp án A. $0,6=3\lambda/2$, suy ra $\lambda=0,4$ m.
}
\end{ex}

% ===== Câu 121 =====
% ID: L11C2B14-31-TL
% Mức: TH
\dangbt{Tổng hợp các đại lượng đặc trưng của sóng}
\begin{ex}
% ID: L11C2B14-31-TL
% Mức: TH
Trong dạng “Tổng hợp các đại lượng đặc trưng của sóng”, Viết biểu thức độ lệch pha của hai điểm trên phương truyền sóng và nêu ý nghĩa.
\choiceTF

\loigiai{
$\Delta\varphi=2\pi\Delta x/\lambda$. Công thức cho biết quan hệ pha phụ thuộc khoảng cách theo phương truyền sóng.
}
\end{ex}

% ===== Câu 122 =====
% ID: L11C2B14-31-TLN
% Mức: TH
\dangbt{Vận dụng tổng hợp kiến thức sóng trong thực tế}
\begin{ex}
% ID: L11C2B14-31-TLN
% Mức: TH
Trong dạng “Vận dụng tổng hợp kiến thức sóng trong thực tế”, Sóng có chu kì $0,25\,\text{s}$. Tính tần số theo Hz.
\shortans{4}
\loigiai{
$f=1/T=1/0,25=4$ Hz.
}
\end{ex}

% ===== Câu 123 =====
% ID: L11C2B14-31-TN
% Mức: NB
\dangbt{Chưa có dạng}
\begin{ex}
% ID: L11C2B14-31-TN
% Mức: NB
Khi có sóng ngang truyền qua, các phần tử vật chất của môi trường dao động $\nguon{ly11 kntt.pdf}$
\choice
{theo phương song song với phương truyền sóng.}
{\True theo phương vuông góc với phương truyền sóng.}
{cùng pha với nhau.}
{với các tần số khác nhau.}
\loigiai{
Chọn B.

Khi có sóng ngang truyền qua, các phần tử vật chất của môi trường dao động theo phương vuông góc với phương truyền sóng.
}
\end{ex}

% ===== Câu 124 =====
% ID: L11C2B14-32-TL
% Mức: VD
\dangbt{Tổng hợp các đại lượng đặc trưng của sóng}
\begin{ex}
% ID: L11C2B14-32-TL
% Mức: VD
Trong dạng “Tổng hợp các đại lượng đặc trưng của sóng”, Hai nguồn cùng pha có $\lambda=3$ cm; tại $M$, $d_1=14$ cm, $d_2=21,5$ cm. Xác định trạng thái giao thoa.
\choiceTF

\loigiai{
$\Delta d=7,5\,\mathrm{cm}=2,5\lambda$, nên $M$ là cực tiểu.
}
\end{ex}

% ===== Câu 125 =====
% ID: L11C2B14-32-TLN
% Mức: TH
\dangbt{Vận dụng tổng hợp kiến thức sóng trong thực tế}
\begin{ex}
% ID: L11C2B14-32-TLN
% Mức: TH
Trong dạng “Vận dụng tổng hợp kiến thức sóng trong thực tế”, Sóng có $v=24$ m/s và $f=8$ Hz. Tính bước sóng theo m.
\shortans{3}
\loigiai{
$\lambda=v/f=24/8=3$ m.
}
\end{ex}

% ===== Câu 126 =====
% ID: L11C2B14-32-TN
% Mức: NB
\dangbt{Chưa có dạng}
\begin{ex}
% ID: L11C2B14-32-TN
% Mức: NB
Trường hợp nào sau đây là một ví dụ về sóng dọc? $\nguon{ly11 kntt.pdf}$
\choice
{Ánh sáng truyền trong không khí.}
{Sóng nước trên mặt hồ.}
{\True Sóng âm lan truyền trong không khí.}
{Sóng truyền một trên sợi dây.}
\loigiai{
Chọn C.

Sóng âm là sóng dọc.
}
\end{ex}

% ===== Câu 127 =====
% ID: L11C2B14-33-TL
% Mức: VD
\dangbt{Tổng hợp các đại lượng đặc trưng của sóng}
\begin{bt}
% ID: L11C2B14-33-TL
% Mức: VD
Trong dạng “Tổng hợp các đại lượng đặc trưng của sóng”, Dây hai đầu cố định dài $1,2\,\text{m}$ có tốc độ sóng $48\,\text{m}$ /s. Tính tần số họa âm thứ ba.
\loigiai{
$f_1=v/(2L)=48/2,4=20$ Hz; $f_3=3f_1=60$ Hz.
}
\end{bt}

% ===== Câu 128 =====
% ID: L11C2B14-33-TLN
% Mức: VD
\dangbt{Vận dụng tổng hợp kiến thức sóng trong thực tế}
\begin{ex}
% ID: L11C2B14-33-TLN
% Mức: VD
Trong dạng “Vận dụng tổng hợp kiến thức sóng trong thực tế”, Hai điểm cách nhau $1\,\text{m}$ trên phương truyền sóng, $\lambda=4$ m. Tính độ lệch pha theo radian, lấy kết quả theo $\pi$.
\shortans{0,5}
\loigiai{
$\Delta\varphi=2\pi\Delta x/\lambda=2\pi/4=0,5\pi$ rad.
}
\end{ex}

% ===== Câu 129 =====
% ID: L11C2B14-33-TN
% Mức: NB
\dangbt{Chưa có dạng}
\begin{ex}
% ID: L11C2B14-33-TN
% Mức: NB
Tất cả các sóng điện từ đều có cùng $\nguon{ly11 kntt.pdf}$
\choice
{tốc độ khi truyền trong một môi trường nhất định.}
{tần số khi truyền trong môi trường chân không.}
{chu kì khi truyền trong một môi trường nhất định.}
{\True tốc độ khi truyền trong chân không.}
\loigiai{
Chọn D.

Khi truyền trong chân không các sóng có tốc độ như nhau.
}
\end{ex}

% ===== Câu 130 =====
% ID: L11C2B14-34-TL
% Mức: VD
\dangbt{Tổng hợp các đại lượng đặc trưng của sóng}
\begin{ex}
% ID: L11C2B14-34-TL
% Mức: VD
Trong dạng “Tổng hợp các đại lượng đặc trưng của sóng”, Giải thích sự thay đổi bước sóng khi sóng truyền sang môi trường khác.
\choiceTF

\loigiai{
Tần số do nguồn quyết định nên không đổi; tốc độ phụ thuộc môi trường, vì $\lambda=v/f$ nên bước sóng thay đổi theo tốc độ.
}
\end{ex}

% ===== Câu 131 =====
% ID: L11C2B14-34-TLN
% Mức: VD
\dangbt{Vận dụng tổng hợp kiến thức sóng trong thực tế}
\begin{ex}
% ID: L11C2B14-34-TLN
% Mức: VD
Trong dạng “Vận dụng tổng hợp kiến thức sóng trong thực tế”, Cường độ âm tăng 100 lần. Mức cường độ âm tăng bao nhiêu dB?
\shortans{20}
\loigiai{
$\Delta L=10\log100=20$ dB.
}
\end{ex}

% ===== Câu 132 =====
% ID: L11C2B14-34-TN
% Mức: NB
\dangbt{Chưa có dạng}
\begin{ex}
% ID: L11C2B14-34-TN
% Mức: NB
Hiện tượng giao thoa sóng xảy ra khi có $\nguon{ly11 kntt.pdf}$
\choice
{hai sóng chuyển động ngược chiều giao nhau.}
{\True hai sóng xuất phát từ hai tâm dao động cùng tần số, cùng pha giao nhau.}
{hai sóng dao động cùng phương, cùng pha giao nhau.}
{hai sóng xuất phát từ hai nguồn dao động cùng tần số.}
\loigiai{
Chọn B.

Hiện tượng giao thoa sóng xảy ra khi có hai sóng xuất phát từ hai tâm dao động cùng tần số, cùng pha giao nhau.
}
\end{ex}

% ===== Câu 133 =====
% ID: L11C2B14-35-TL
% Mức: VDC
\dangbt{Tổng hợp các đại lượng đặc trưng của sóng}
\begin{ex}
% ID: L11C2B14-35-TL
% Mức: VDC
Trong dạng “Tổng hợp các đại lượng đặc trưng của sóng”, Tại sao đứng xa nguồn âm thì âm nghe nhỏ hơn?
\choiceTF

\loigiai{
Năng lượng sóng phân bố trên diện tích ngày càng lớn; với nguồn điểm, cường độ giảm theo $1/r^2$, nên âm nghe nhỏ hơn.
}
\end{ex}

% ===== Câu 134 =====
% ID: L11C2B14-35-TLN
% Mức: VD
\dangbt{Vận dụng tổng hợp kiến thức sóng trong thực tế}
\begin{ex}
% ID: L11C2B14-35-TLN
% Mức: VD
Trong dạng “Vận dụng tổng hợp kiến thức sóng trong thực tế”, Dây hai đầu cố định dài $1\,\text{m}$ có bước sóng $0,5\,\text{m}$. Số bụng là bao nhiêu?
\shortans{4}
\loigiai{
$k=2L/\lambda=2/0,5=4$ bụng.
}
\end{ex}

% ===== Câu 135 =====
% ID: L11C2B14-35-TN
% Mức: NB
\dangbt{Chưa có dạng}
\begin{ex}
% ID: L11C2B14-35-TN
% Mức: NB
D. Sóng được tạo thành do sự giao thoa giữa sóng tới và sóng phản xạ. Một sóng vô tuyến được phát ra từ một đài phát thanh có bước sóng $3\,\text{m}$. Coi rằng tốc độ của sóng vô tuyến truyền trong không khí là $3.108 \mathrm{m/s}$, tần số của sóng này là $\nguon{ly11 kntt.pdf}$
\choice
{$10^{-8} \mathrm{Hz}.$}
{$9\cdot 10^{-8} \mathrm{Hz}.$}
{\True $108 \mathrm{Hz}$.}
{$9 \mathrm{Hz}$}
\loigiai{
Chọn C.

$f = \dfrac{v}{\lambda} = \dfrac{3\cdot 10^{8}}{3} = 10^{8}Hz$
}
\end{ex}

% ===== Câu 136 =====
% ID: L11C2B14-36-TL
% Mức: TH
\dangbt{Vận dụng tổng hợp kiến thức sóng trong thực tế}
\begin{bt}
% ID: L11C2B14-36-TL
% Mức: TH
Trong dạng “Vận dụng tổng hợp kiến thức sóng trong thực tế”, Một sóng truyền với tốc độ $20\,\text{m}$ /s và chu kì $0,1\,\text{s}$. Tính bước sóng.
\loigiai{
$f=1/T=10$ Hz; $\lambda=v/f=20/10=2$ m.
}
\end{bt}

% ===== Câu 137 =====
% ID: L11C2B14-36-TLN
% Mức: VDC
\begin{ex}
% ID: L11C2B14-36-TLN
% Mức: VDC
Trong dạng “Vận dụng tổng hợp kiến thức sóng trong thực tế”, Nguồn điểm có cường độ âm 8 đơn vị tại r. Ở khoảng cách 2r, cường độ theo cùng đơn vị là bao nhiêu?
\shortans{2}
\loigiai{
$I_2=I_1(r/2r)^2=8/4=2$.
}
\end{ex}

% ===== Câu 138 =====
% ID: L11C2B14-36-TN
% Mức: NB
\dangbt{Chưa có dạng}
\begin{ex}
% ID: L11C2B14-36-TN
% Mức: NB
Người ta tạo ra sóng dừng trên một sợi dây căng giữa hai điểm cố định. Hai tần số gần nhau nhất cùng tạo ra sóng dừng trên dây là $150 \mathrm{Hz}$ và $200 \mathrm{Hz}$. Tần số nhỏ nhất tạo ra sóng dừng trên dây đó là $\nguon{ly11 kntt.pdf}$
\choice
{\True $50 \mathrm{Hz}$.}
{$75 \mathrm{Hz}$.}
{$100 \mathrm{Hz}$.}
{$125\,\mathrm{H}$}
\loigiai{
Chọn A.

Vì sợi dây có hai đầu cố định nên: f_(min) = f_(k + 1) - f_(k)= 200 - 150 = $50\,\mathrm{Hz}$
}
\end{ex}

% ===== Câu 139 =====
% ID: L11C2B14-37-TL
% Mức: TH
\dangbt{Vận dụng tổng hợp kiến thức sóng trong thực tế}
\begin{ex}
% ID: L11C2B14-37-TL
% Mức: TH
Trong dạng “Vận dụng tổng hợp kiến thức sóng trong thực tế”, Viết biểu thức độ lệch pha của hai điểm trên phương truyền sóng và nêu ý nghĩa.
\choiceTF

\loigiai{
$\Delta\varphi=2\pi\Delta x/\lambda$. Công thức cho biết quan hệ pha phụ thuộc khoảng cách theo phương truyền sóng.
}
\end{ex}

% ===== Câu 140 =====
% ID: L11C2B14-37-TN
% Mức: NB
\dangbt{Chưa có dạng}
\begin{ex}
% ID: L11C2B14-37-TN
% Mức: NB
Trong thí nghiệm Young về giao thoa với ánh sáng đơn sắc có bước sóng $\lambda$. Biết khoảng cách giữa hai khe là $1 \mathrm{mm}$. Tại điểm $M$ cách vân sáng trung tâm $1,2 \mathrm{mm}$ trên màn quan sát là vị trí vân sáng bậc 4. Nếu dịch màn ra xa thêm một đoạn $25 \mathrm{cm}$ theo phương vuông góc với mặt phẳng hai khe thi tại $M$ là vị trí vân sáng bậc 3. Bước sóng $\lambda$ dùng trong thí nghiệm là $\nguon{ly11 kntt.pdf}$
\choice
{\True 0,4 $\mu$ m.}
{0,5 $\mu$ m.}
{0,6 $\mu$ m.}
{0,64 $\mu$}
\loigiai{
Chọn A.
}
\end{ex}

% ===== Câu 141 =====
% ID: L11C2B14-38-TL
% Mức: VD
\dangbt{Vận dụng tổng hợp kiến thức sóng trong thực tế}
\begin{ex}
% ID: L11C2B14-38-TL
% Mức: VD
Trong dạng “Vận dụng tổng hợp kiến thức sóng trong thực tế”, Hai nguồn cùng pha có $\lambda=3$ cm; tại $M$, $d_1=14$ cm, $d_2=21,5$ cm. Xác định trạng thái giao thoa.
\choiceTF

\loigiai{
$\Delta d=7,5\,\mathrm{cm}=2,5\lambda$, nên $M$ là cực tiểu.
}
\end{ex}

% ===== Câu 142 =====
% ID: L11C2B14-38-TN
% Mức: NB
\dangbt{Phương trình sóng và độ lệch pha}
\begin{ex}
% ID: L11C2B14-38-TN
% Mức: NB
Trong dạng “Phương trình sóng và độ lệch pha”, Một sóng có $v=30$ m/s, $f=5$ Hz. Bước sóng là
\choice
{\True $6\,\text{m}$}
{$150\,\text{m}$}
{$0,167\,\text{m}$}
{$35\,\text{m}$}
\loigiai{
Đáp án A. $\lambda=v/f=30/5=6$ m.
}
\end{ex}

% ===== Câu 143 =====
% ID: L11C2B14-39-TL
% Mức: VD
\dangbt{Vận dụng tổng hợp kiến thức sóng trong thực tế}
\begin{bt}
% ID: L11C2B14-39-TL
% Mức: VD
Trong dạng “Vận dụng tổng hợp kiến thức sóng trong thực tế”, Dây hai đầu cố định dài $1,2\,\text{m}$ có tốc độ sóng $48\,\text{m}$ /s. Tính tần số họa âm thứ ba.
\loigiai{
$f_1=v/(2L)=48/2,4=20$ Hz; $f_3=3f_1=60$ Hz.
}
\end{bt}

% ===== Câu 144 =====
% ID: L11C2B14-39-TN
% Mức: NB
\dangbt{Phương trình sóng và độ lệch pha}
\begin{ex}
% ID: L11C2B14-39-TN
% Mức: NB
Trong dạng “Phương trình sóng và độ lệch pha”, Độ lệch pha giữa hai điểm cách nhau $\Delta x$ trên phương truyền sóng là
\choice
{\True $2\pi\Delta x/\lambda$}
{$\pi\Delta x/\lambda$}
{$2\pi\lambda/\Delta x$}
{$\Delta x/(2\pi\lambda)$}
\loigiai{
Đáp án A. Theo phương trình sóng, $|\Delta\varphi|=2\pi\Delta x/\lambda$.
}
\end{ex}

% ===== Câu 145 =====
% ID: L11C2B14-40-TL
% Mức: VD
\dangbt{Vận dụng tổng hợp kiến thức sóng trong thực tế}
\begin{ex}
% ID: L11C2B14-40-TL
% Mức: VD
Trong dạng “Vận dụng tổng hợp kiến thức sóng trong thực tế”, Giải thích sự thay đổi bước sóng khi sóng truyền sang môi trường khác.
\choiceTF

\loigiai{
Tần số do nguồn quyết định nên không đổi; tốc độ phụ thuộc môi trường, vì $\lambda=v/f$ nên bước sóng thay đổi theo tốc độ.
}
\end{ex}

% ===== Câu 146 =====
% ID: L11C2B14-40-TN
% Mức: TH
\dangbt{Phương trình sóng và độ lệch pha}
\begin{ex}
% ID: L11C2B14-40-TN
% Mức: TH
Trong dạng “Phương trình sóng và độ lệch pha”, Hai điểm cách nhau $\lambda/4$ trên phương truyền sóng lệch pha
\choice
{\True $\pi/2$}
{$\pi$}
{$2\pi$}
{$\pi/4$}
\loigiai{
Đáp án A. $\Delta\varphi=2\pi(\lambda/4)/\lambda=\pi/2$.
}
\end{ex}

% ===== Câu 147 =====
% ID: L11C2B14-41-TL
% Mức: VDC
\dangbt{Vận dụng tổng hợp kiến thức sóng trong thực tế}
\begin{ex}
% ID: L11C2B14-41-TL
% Mức: VDC
Trong dạng “Vận dụng tổng hợp kiến thức sóng trong thực tế”, Tại sao đứng xa nguồn âm thì âm nghe nhỏ hơn?
\choiceTF

\loigiai{
Năng lượng sóng phân bố trên diện tích ngày càng lớn; với nguồn điểm, cường độ giảm theo $1/r^2$, nên âm nghe nhỏ hơn.
}
\end{ex}

% ===== Câu 148 =====
% ID: L11C2B14-41-TN
% Mức: TH
\dangbt{Phương trình sóng và độ lệch pha}
\begin{ex}
% ID: L11C2B14-41-TN
% Mức: TH
Trong dạng “Phương trình sóng và độ lệch pha”, Hai nguồn cùng pha, cực đại giao thoa thỏa
\choice
{\True $\Delta d=k\lambda$}
{$\Delta d=(k+1/2)\lambda$}
{$d_1+d_2=k\lambda$}
{$\Delta d=\lambda/4$}
\loigiai{
Đáp án A. Điều kiện cực đại của hai nguồn cùng pha là $\Delta d=k\lambda$.
}
\end{ex}

% ===== Câu 149 =====
% ID: L11C2B14-42-TN
% Mức: TH
\begin{ex}
% ID: L11C2B14-42-TN
% Mức: TH
Trong dạng “Phương trình sóng và độ lệch pha”, Hai nút sóng dừng liên tiếp cách nhau
\choice
{\True $\lambda/2$}
{$\lambda$}
{$\lambda/4$}
{$2\lambda$}
\loigiai{
Đáp án A. Khoảng cách hai nút liên tiếp là nửa bước sóng.
}
\end{ex}

% ===== Câu 150 =====
% ID: L11C2B14-43-TN
% Mức: VD
\begin{ex}
% ID: L11C2B14-43-TN
% Mức: VD
Trong dạng “Phương trình sóng và độ lệch pha”, Âm có tần số càng lớn thì cảm giác âm càng
\choice
{\True cao}
{to}
{trầm}
{nhỏ}
\loigiai{
Đáp án A. Độ cao của âm phụ thuộc chủ yếu vào tần số.
}
\end{ex}

% ===== Câu 151 =====
% ID: L11C2B14-44-TN
% Mức: VD
\begin{ex}
% ID: L11C2B14-44-TN
% Mức: VD
Trong dạng “Phương trình sóng và độ lệch pha”, Mức cường độ âm được tính bởi
\choice
{\True $L=10\log(I/I_0)$}
{$L=I/I_0$}
{$L=20I/I_0$}
{$L=10I_0/I$}
\loigiai{
Đáp án A. Theo định nghĩa, $L=10\log_{10}(I/I_0)$, đơn vị dB.
}
\end{ex}

% ===== Câu 152 =====
% ID: L11C2B14-45-TN
% Mức: VD
\begin{ex}
% ID: L11C2B14-45-TN
% Mức: VD
Trong dạng “Phương trình sóng và độ lệch pha”, Khi khoảng cách đến nguồn điểm tăng gấp đôi, cường độ âm trong môi trường không hấp thụ
\choice
{\True giảm 4 lần}
{giảm 2 lần}
{tăng 4 lần}
{không đổi}
\loigiai{
Đáp án A. Với nguồn điểm $I\sim1/r^2$, nên r tăng 2 lần thì $I$ giảm 4 lần.
}
\end{ex}

% ===== Câu 153 =====
% ID: L11C2B14-46-TN
% Mức: VDC
\begin{ex}
% ID: L11C2B14-46-TN
% Mức: VDC
Trong dạng “Phương trình sóng và độ lệch pha”, Sóng truyền từ môi trường này sang môi trường khác, đại lượng không đổi là
\choice
{\True tần số}
{bước sóng}
{tốc độ}
{biên độ}
\loigiai{
Đáp án A. Tần số do nguồn quyết định và không đổi khi qua mặt phân cách.
}
\end{ex}

% ===== Câu 154 =====
% ID: L11C2B14-47-TN
% Mức: VDC
\begin{ex}
% ID: L11C2B14-47-TN
% Mức: VDC
Trong dạng “Phương trình sóng và độ lệch pha”, Một dây hai đầu cố định dài $0,6\,\text{m}$ có 3 bụng. Bước sóng là
\choice
{\True $0,4\,\text{m}$}
{$0,2\,\text{m}$}
{$1,8\,\text{m}$}
{$0,6\,\text{m}$}
\loigiai{
Đáp án A. $0,6=3\lambda/2$, suy ra $\lambda=0,4$ m.
}
\end{ex}

% ===== Câu 155 =====
% ID: L11C2B14-48-TN
% Mức: NB
\dangbt{Tổng hợp các đại lượng đặc trưng của sóng}
\begin{ex}
% ID: L11C2B14-48-TN
% Mức: NB
Trong dạng “Tổng hợp các đại lượng đặc trưng của sóng”, Một sóng có $v=30$ m/s, $f=5$ Hz. Bước sóng là
\choice
{\True $6\,\text{m}$}
{$150\,\text{m}$}
{$0,167\,\text{m}$}
{$35\,\text{m}$}
\loigiai{
Đáp án A. $\lambda=v/f=30/5=6$ m.
}
\end{ex}

% ===== Câu 156 =====
% ID: L11C2B14-49-TN
% Mức: NB
\begin{ex}
% ID: L11C2B14-49-TN
% Mức: NB
Trong dạng “Tổng hợp các đại lượng đặc trưng của sóng”, Độ lệch pha giữa hai điểm cách nhau $\Delta x$ trên phương truyền sóng là
\choice
{\True $2\pi\Delta x/\lambda$}
{$\pi\Delta x/\lambda$}
{$2\pi\lambda/\Delta x$}
{$\Delta x/(2\pi\lambda)$}
\loigiai{
Đáp án A. Theo phương trình sóng, $|\Delta\varphi|=2\pi\Delta x/\lambda$.
}
\end{ex}

% ===== Câu 157 =====
% ID: L11C2B14-50-TN
% Mức: TH
\begin{ex}
% ID: L11C2B14-50-TN
% Mức: TH
Trong dạng “Tổng hợp các đại lượng đặc trưng của sóng”, Hai điểm cách nhau $\lambda/4$ trên phương truyền sóng lệch pha
\choice
{\True $\pi/2$}
{$\pi$}
{$2\pi$}
{$\pi/4$}
\loigiai{
Đáp án A. $\Delta\varphi=2\pi(\lambda/4)/\lambda=\pi/2$.
}
\end{ex}

% ===== Câu 158 =====
% ID: L11C2B14-51-TN
% Mức: TH
\begin{ex}
% ID: L11C2B14-51-TN
% Mức: TH
Trong dạng “Tổng hợp các đại lượng đặc trưng của sóng”, Hai nguồn cùng pha, cực đại giao thoa thỏa
\choice
{\True $\Delta d=k\lambda$}
{$\Delta d=(k+1/2)\lambda$}
{$d_1+d_2=k\lambda$}
{$\Delta d=\lambda/4$}
\loigiai{
Đáp án A. Điều kiện cực đại của hai nguồn cùng pha là $\Delta d=k\lambda$.
}
\end{ex}

% ===== Câu 159 =====
% ID: L11C2B14-52-TN
% Mức: TH
\begin{ex}
% ID: L11C2B14-52-TN
% Mức: TH
Trong dạng “Tổng hợp các đại lượng đặc trưng của sóng”, Hai nút sóng dừng liên tiếp cách nhau
\choice
{\True $\lambda/2$}
{$\lambda$}
{$\lambda/4$}
{$2\lambda$}
\loigiai{
Đáp án A. Khoảng cách hai nút liên tiếp là nửa bước sóng.
}
\end{ex}

% ===== Câu 160 =====
% ID: L11C2B14-53-TN
% Mức: VD
\begin{ex}
% ID: L11C2B14-53-TN
% Mức: VD
Trong dạng “Tổng hợp các đại lượng đặc trưng của sóng”, Âm có tần số càng lớn thì cảm giác âm càng
\choice
{\True cao}
{to}
{trầm}
{nhỏ}
\loigiai{
Đáp án A. Độ cao của âm phụ thuộc chủ yếu vào tần số.
}
\end{ex}

% ===== Câu 161 =====
% ID: L11C2B14-54-TN
% Mức: VD
\begin{ex}
% ID: L11C2B14-54-TN
% Mức: VD
Trong dạng “Tổng hợp các đại lượng đặc trưng của sóng”, Mức cường độ âm được tính bởi
\choice
{\True $L=10\log(I/I_0)$}
{$L=I/I_0$}
{$L=20I/I_0$}
{$L=10I_0/I$}
\loigiai{
Đáp án A. Theo định nghĩa, $L=10\log_{10}(I/I_0)$, đơn vị dB.
}
\end{ex}

% ===== Câu 162 =====
% ID: L11C2B14-55-TN
% Mức: VD
\begin{ex}
% ID: L11C2B14-55-TN
% Mức: VD
Trong dạng “Tổng hợp các đại lượng đặc trưng của sóng”, Khi khoảng cách đến nguồn điểm tăng gấp đôi, cường độ âm trong môi trường không hấp thụ
\choice
{\True giảm 4 lần}
{giảm 2 lần}
{tăng 4 lần}
{không đổi}
\loigiai{
Đáp án A. Với nguồn điểm $I\sim1/r^2$, nên r tăng 2 lần thì $I$ giảm 4 lần.
}
\end{ex}

% ===== Câu 163 =====
% ID: L11C2B14-56-TN
% Mức: VDC
\begin{ex}
% ID: L11C2B14-56-TN
% Mức: VDC
Trong dạng “Tổng hợp các đại lượng đặc trưng của sóng”, Sóng truyền từ môi trường này sang môi trường khác, đại lượng không đổi là
\choice
{\True tần số}
{bước sóng}
{tốc độ}
{biên độ}
\loigiai{
Đáp án A. Tần số do nguồn quyết định và không đổi khi qua mặt phân cách.
}
\end{ex}

% ===== Câu 164 =====
% ID: L11C2B14-57-TN
% Mức: VDC
\begin{ex}
% ID: L11C2B14-57-TN
% Mức: VDC
Trong dạng “Tổng hợp các đại lượng đặc trưng của sóng”, Một dây hai đầu cố định dài $0,6\,\text{m}$ có 3 bụng. Bước sóng là
\choice
{\True $0,4\,\text{m}$}
{$0,2\,\text{m}$}
{$1,8\,\text{m}$}
{$0,6\,\text{m}$}
\loigiai{
Đáp án A. $0,6=3\lambda/2$, suy ra $\lambda=0,4$ m.
}
\end{ex}

% ===== Câu 165 =====
% ID: L11C2B14-58-TN
% Mức: NB
\dangbt{Vận dụng tổng hợp kiến thức sóng trong thực tế}
\begin{ex}
% ID: L11C2B14-58-TN
% Mức: NB
Trong dạng “Vận dụng tổng hợp kiến thức sóng trong thực tế”, Một sóng có $v=30$ m/s, $f=5$ Hz. Bước sóng là
\choice
{\True $6\,\text{m}$}
{$150\,\text{m}$}
{$0,167\,\text{m}$}
{$35\,\text{m}$}
\loigiai{
Đáp án A. $\lambda=v/f=30/5=6$ m.
}
\end{ex}

% ===== Câu 166 =====
% ID: L11C2B14-59-TN
% Mức: NB
\begin{ex}
% ID: L11C2B14-59-TN
% Mức: NB
Trong dạng “Vận dụng tổng hợp kiến thức sóng trong thực tế”, Độ lệch pha giữa hai điểm cách nhau $\Delta x$ trên phương truyền sóng là
\choice
{\True $2\pi\Delta x/\lambda$}
{$\pi\Delta x/\lambda$}
{$2\pi\lambda/\Delta x$}
{$\Delta x/(2\pi\lambda)$}
\loigiai{
Đáp án A. Theo phương trình sóng, $|\Delta\varphi|=2\pi\Delta x/\lambda$.
}
\end{ex}

% ===== Câu 167 =====
% ID: L11C2B14-60-TN
% Mức: TH
\begin{ex}
% ID: L11C2B14-60-TN
% Mức: TH
Trong dạng “Vận dụng tổng hợp kiến thức sóng trong thực tế”, Hai điểm cách nhau $\lambda/4$ trên phương truyền sóng lệch pha
\choice
{\True $\pi/2$}
{$\pi$}
{$2\pi$}
{$\pi/4$}
\loigiai{
Đáp án A. $\Delta\varphi=2\pi(\lambda/4)/\lambda=\pi/2$.
}
\end{ex}

% ===== Câu 168 =====
% ID: L11C2B14-61-TN
% Mức: TH
\begin{ex}
% ID: L11C2B14-61-TN
% Mức: TH
Trong dạng “Vận dụng tổng hợp kiến thức sóng trong thực tế”, Hai nguồn cùng pha, cực đại giao thoa thỏa
\choice
{\True $\Delta d=k\lambda$}
{$\Delta d=(k+1/2)\lambda$}
{$d_1+d_2=k\lambda$}
{$\Delta d=\lambda/4$}
\loigiai{
Đáp án A. Điều kiện cực đại của hai nguồn cùng pha là $\Delta d=k\lambda$.
}
\end{ex}

% ===== Câu 169 =====
% ID: L11C2B14-62-TN
% Mức: TH
\begin{ex}
% ID: L11C2B14-62-TN
% Mức: TH
Trong dạng “Vận dụng tổng hợp kiến thức sóng trong thực tế”, Hai nút sóng dừng liên tiếp cách nhau
\choice
{\True $\lambda/2$}
{$\lambda$}
{$\lambda/4$}
{$2\lambda$}
\loigiai{
Đáp án A. Khoảng cách hai nút liên tiếp là nửa bước sóng.
}
\end{ex}

% ===== Câu 170 =====
% ID: L11C2B14-63-TN
% Mức: VD
\begin{ex}
% ID: L11C2B14-63-TN
% Mức: VD
Trong dạng “Vận dụng tổng hợp kiến thức sóng trong thực tế”, Âm có tần số càng lớn thì cảm giác âm càng
\choice
{\True cao}
{to}
{trầm}
{nhỏ}
\loigiai{
Đáp án A. Độ cao của âm phụ thuộc chủ yếu vào tần số.
}
\end{ex}

% ===== Câu 171 =====
% ID: L11C2B14-64-TN
% Mức: VD
\begin{ex}
% ID: L11C2B14-64-TN
% Mức: VD
Trong dạng “Vận dụng tổng hợp kiến thức sóng trong thực tế”, Mức cường độ âm được tính bởi
\choice
{\True $L=10\log(I/I_0)$}
{$L=I/I_0$}
{$L=20I/I_0$}
{$L=10I_0/I$}
\loigiai{
Đáp án A. Theo định nghĩa, $L=10\log_{10}(I/I_0)$, đơn vị dB.
}
\end{ex}

% ===== Câu 172 =====
% ID: L11C2B14-65-TN
% Mức: VD
\begin{ex}
% ID: L11C2B14-65-TN
% Mức: VD
Trong dạng “Vận dụng tổng hợp kiến thức sóng trong thực tế”, Khi khoảng cách đến nguồn điểm tăng gấp đôi, cường độ âm trong môi trường không hấp thụ
\choice
{\True giảm 4 lần}
{giảm 2 lần}
{tăng 4 lần}
{không đổi}
\loigiai{
Đáp án A. Với nguồn điểm $I\sim1/r^2$, nên r tăng 2 lần thì $I$ giảm 4 lần.
}
\end{ex}

% ===== Câu 173 =====
% ID: L11C2B14-66-TN
% Mức: VDC
\begin{ex}
% ID: L11C2B14-66-TN
% Mức: VDC
Trong dạng “Vận dụng tổng hợp kiến thức sóng trong thực tế”, Sóng truyền từ môi trường này sang môi trường khác, đại lượng không đổi là
\choice
{\True tần số}
{bước sóng}
{tốc độ}
{biên độ}
\loigiai{
Đáp án A. Tần số do nguồn quyết định và không đổi khi qua mặt phân cách.
}
\end{ex}

% ===== Câu 174 =====
% ID: L11C2B14-67-TN
% Mức: VDC
\begin{ex}
% ID: L11C2B14-67-TN
% Mức: VDC
Trong dạng “Vận dụng tổng hợp kiến thức sóng trong thực tế”, Một dây hai đầu cố định dài $0,6\,\text{m}$ có 3 bụng. Bước sóng là
\choice
{\True $0,4\,\text{m}$}
{$0,2\,\text{m}$}
{$1,8\,\text{m}$}
{$0,6\,\text{m}$}
\loigiai{
Đáp án A. $0,6=3\lambda/2$, suy ra $\lambda=0,4$ m.
}
\end{ex}
